\documentclass[times]{itmo-student-thesis}

%% Опции пакета:
%% - times          — шрифт Times New Roman (сборка через xelatex)
%% В PDF ВКР не включаем задание/аннотацию/титульные листы класса: только
%% оглавление и основной текст (титул и задание оформляются отдельно по ИТМО).

%% Оформление шрифтов по ГОСТ 7.32: основной текст --- 14 pt (класс
%% itmo-student-thesis), полуторный интервал --- в классе (\onehalfspacing).
%% Times New Roman: при установленном семействе в ОС можно заменить строки
%% ниже на \setmainfont{Times New Roman}[Ligatures={TeX,Historic}] и т.д.
\setmainfont[Ligatures={TeX,Historic}]{Liberation Serif}
\setsansfont{Liberation Sans}
\setmonofont{Liberation Mono}

%% Поля страницы (ГОСТ 7.32-2017, типовые для текстового документа А4):
%% левое $\ge$30~мм под переплёт, прочие --- по стандарту; при расхождении
%% с локальными указаниями ИТМО следует руководствоваться методичкой кафедры.
\geometry{left=30mm,right=15mm,top=20mm,bottom=20mm}

%% extreport: в \begin{titlepage} задаётся \setcounter{page}{1}; при оглавлении
%% первым это сбивает нумерацию. Оставляем \thispagestyle{empty}, без сброса счётчика.
\makeatletter
\renewenvironment{titlepage}{%
  \if@twocolumn
    \@restonecoltrue\onecolumn
  \else
    \@restonecolfalse\newpage
  \fi
  \thispagestyle{empty}%
}{%
  \if@restonecol\twocolumn\else\newpage\fi
}
\makeatother

%% В PDF только оглавление + текст: \maketitle класса ИТМО без титулов/задания/аннотации.
\makeatletter
\renewcommand{\maketitle}[1]{%
    \ifstrequal{#1}{Магистр}{%
        \providecommand{\thespecialization@ru}{Технологии разработки программного обеспечения}%
        \providecommand{\thespecialty@ru}{01.04.02~Прикладная математика и информатика}%
        \providecommand{\thespecialization@en}{Technologies of software engineering}%
        \providecommand{\thespecialty@en}{01.04.02~Applied mathematics and informatics}%
        \providecommand{\@strings@DegreeName@ru}{Магистр}%
        \providecommand{\@strings@DegreeName@en}{Master}%
    }{%
        \ifstrequal{#1}{Бакалавр}{%
            \providecommand{\thespecialization@ru}{Информатика и программирование}%
            \providecommand{\thespecialty@ru}{01.03.02~Прикладная математика и информатика}%
            \providecommand{\thespecialization@en}{Computer science and programming}%
            \providecommand{\thespecialty@en}{01.03.02~Applied mathematics and informatics}%
            \providecommand{\@strings@DegreeName@ru}{Бакалавр}%
            \providecommand{\@strings@DegreeName@en}{Bachelor}%
        }{%
            \PackageError{thesis}{Unknown maketitle argument: #1, expected Магистр or Бакалавр}{}%
        }%
    }%
    \providecommand{\TheOnlyTruePageStyle}{fancy}%
    \pagestyle{\TheOnlyTruePageStyle}%
    \renewcommand{\headrulewidth}{0pt}%
    \renewcommand{\footrulewidth}{0pt}%
    \lhead{}%
    \chead{\thepage}%
    \rhead{}%
    \lfoot{}%
    \cfoot{}%
    \rfoot{}%
    \addtolength{\voffset}{3mm}%
    \addtolength{\headsep}{-3mm}%
    \pretocmd{\tableofcontents}{\begingroup\banhyphens}{}{}%
    \apptocmd{\tableofcontents}{\endgroup}{}{}%
}
\makeatother

\usepackage{icomma}       % «умная» запятая в формулах
%% tabularx, listings, xcolor уже загружены в itmo-student-thesis.cls
\usepackage{booktabs}     % \toprule / \midrule / \bottomrule
\usepackage{graphicx}      % для включения диаграмм из PlantUML
\usepackage{float}         % опция [H] — рисунок на месте, без переноса на отдельную страницу
\usepackage{needspace}    % \Needspace — не отрывать блок от начала страницы (таблицы с текстом)
\usepackage{tikz}
\usetikzlibrary{arrows.meta, positioning, shapes.geometric, fit, calc}

%% Настройка listings для Python (cls уже задал базовые параметры)
\definecolor{codegray}{rgb}{0.5,0.5,0.5}
\definecolor{codegreen}{rgb}{0,0.5,0}
\definecolor{codeblue}{rgb}{0,0,0.6}
\lstset{
    language=Python,
    basicstyle=\small\ttfamily,
    keywordstyle=\color{codeblue}\bfseries,
    commentstyle=\color{codegreen},
    stringstyle=\color{codegray},
    numbers=left,
    numberstyle=\tiny\color{codegray},
    breaklines=true,
    frame=single,
    captionpos=b
}

%% Подписи к листингам: снизу по центру (как у рисунков в шаблоне ИТМО).
\AtBeginDocument{%
    \floatsetup[lstlisting]{style=plain, capposition=bottom}%
    \captionsetup[lstlisting]{justification=centering, position=bottom}%
}

\addbibresource{thesis.bib}

%% Скрытые проходы \printbibliography внутри \printmainbibliography (класс ИТМО)
%% для подсчёта источников разбивают страницы --- между списком литературы
%% и приложениями появляется пустой лист. Выполняем подсчёт в \vbox без вывода.
\makeatletter
\renewcommand{\@countsourceswhere}[2]{%
  \setcounter{@bibliocounter}{0}%
  \setbox\@tempboxa=\vbox\bgroup
    \begingroup\@killhyperrefissue
    \printbibliography[env=@fillbibliocounter,heading=reallynone,#1]%
    \endgroup
  \egroup
  \setcounter{#2}{\arabic{@bibliocounter}}}
\makeatother

%%=============================================================
%% ДОКУМЕНТ
%%=============================================================

\begin{document}

%%=============================================================
%% ДАННЫЕ РАБОТЫ (должны быть внутри \begin{document})
%%=============================================================

%% TODO: укажите свою учебную группу
\studygroup{M3434}

\title{Разработка инструмента анализа информационных контекстов,
       направленного на ускорение и повышение качества текстовой разметки}

%% TODO: укажите своё отчество
\author{Новицкий Илья Андреевич}{Новицкий И.\,А.}

%% TODO: уточните должность/степень и место работы руководителя
\supervisor{Ефимова Валерия Александровна}{Ефимова В.\,А.}
           {доцент, кандидат технических наук}
           {Университет ИТМО}

\publishyear{2026}

%% Даты (можно убрать — поля будут заполнены вручную)
\startdate{01}{сентября}{2025}
\finishdate{31}{мая}{2026}
\defencedate{15}{июня}{2026}

%% TODO: уточните должность/степень консультанта
\addconsultant{Ильичёв А.\,С.}{руководитель бригады прикладных генеративных сценариев, магистр, ООО «Яндекс.Технологии»}

%% TODO: укажите фамилию и.о. секретаря ГЭК
\secretary{Ефимова В.\,А.}

%%=============================================================
%% ЗАДАНИЕ
%%=============================================================

\technicalspec{%
Улучшить систему разметки фрагментов ответов LLM: ускорить процесс в~два раза
до~3--4~минут на задание (вдвое сократив время поиска контекстов), повысить
точность до~F-меры выше~$0{,}85$ (с уровня~$\approx 0{,}60$), обеспечить
латентность ответа специализированной модели~$\le 30$~секунд на 1~фрагмент
из задания (вместо~2--5~минут на 1~фрагмент из задания у внешней LLM)
за счёт спекулятивного декодирования
(Eagle) и квантизации в~FP8, снизить долю галлюцинаций (ориентир с~$\approx 30$\%)
и поднять удовлетворённость команды выше~70\% (с~$\approx 50$\%).
Инструмент~--- дистиллированная языковая модель в платформе текстовой разметки
ООО «Яндекс.Технологии»: вход~--- информационные контексты и запрос исполнителя,
выход~--- структурированный ответ о степени подтверждённости фактов.%
}

\plannedcontents{%
\begin{enumerate}
    \item Аналитический обзор методов ускоренной разметки и оценки достоверности
          ответов LLM: процесс разметки, источники трудозатрат, существующие
          платформы и LLM-ассистенты; методы дистилляции, дообучения и оценки
          качества; выбор модели-ученика и технологического стека.
    \item Проектирование инструмента анализа информационных контекстов:
          требования к инструменту; архитектура (C4, UML Sequence, конвейеры);
          оркестрация; логическая модель данных; инференс-сервис.
    \item Реализация, тестирование и оптимизация инструмента разметки:
          пайплайны сбора данных и генерации датасета; обучение (Full fine-tune,
          LoRA, CE~RL); метрики regex и LLM-as-a-judge; инференс с оптимизацией
          FP8 и Eagle; регрессия, нагрузка, внедрение; оценка эффекта «до/после».
    \item (В рамках главы~3) Оптимизация инференса: квантизация FP8, спекулятивное
          декодирование, методика замеров и сравнение конфигураций.
\end{enumerate}%
}

\plannedsources{%
\begin{enumerate}
    \item Hinton G., Vinyals O., Dean J. Distilling the Knowledge in a Neural
          Network // arXiv preprint arXiv:1503.02531. 2015.
    \item Hu E.\,J. et al. LoRA: Low-Rank Adaptation of Large Language Models
          // ICLR. 2022.
    \item Zheng L. et al. Judging LLM-as-a-Judge with MT-Bench and Chatbot Arena
          // NeurIPS. 2023.
    \item Dao T. et al. FlashAttention // NeurIPS. 2022.
    \item Rajbhandari S. et al. ZeRO: Memory Optimizations Toward Training
          Trillion Parameter Models // SC20. 2020.
    \item Внутренняя документация платформы разметки ООО «Яндекс.Технологии».
\end{enumerate}%
}

%%=============================================================
%% АННОТАЦИЯ
%%=============================================================

\researchaim{%
Улучшить систему разметки фрагментов: ускорить процесс в~два раза до~3--4~минут
на задание, повысить F-меру выше~$0{,}85$ (с~$\approx 0{,}60$), обеспечить
латентность~$\le 30$~с на 1~фрагмент из задания за счёт Eagle и квантизации FP8, снизить галлюцинации
и поднять удовлетворённость команды выше~70\% при внедрении дистиллированной
модели в платформу разметки ООО «Яндекс.Технологии».%
}

\researchtargets{%
\begin{enumerate}
    \item Подобрать архитектуру и размер генеративной модели-«ученика»
          с балансом скорость/качество и спроектировать процедуру её обучения.
    \item Сформировать синтетические датасеты запросов и ответов
          (distillation from large models), реализовать очистку/фильтрацию
          и контроль распределений.
    \item Реализовать пайплайны оркестрации (статический и динамический граф)
          для подготовки данных, генерации, обучения, оценки и инференса.
    \item Разработать метрики автоматической оценки качества ответов
          (правила/регулярные выражения + LLM-as-a-judge) и провести серию
          экспериментов по гиперпараметрам и оптимизациям обучения/инференса.
\end{enumerate}%
}

\addadvancedsoftware{\texttt{PyTorch}~— фреймворк машинного обучения}
    {Главы~2, 3}
\addadvancedsoftware{Внутренний аналог Apache Airflow (статический и динамический)
    ~— оркестрация рабочих процессов}{Главы~2, 3}
\addadvancedsoftware{\texttt{YTSaurus}~— платформа распределённого хранения данных}
    {Главы~2, 3}
\addadvancedsoftware{Внутренний аналог \texttt{vLLM}~— сервер ускоренного инференса}
    {Глава~3}
\addadvancedsoftware{\texttt{NLTK}~— обработка естественного языка}{Глава~3}

\researchsummary{%
Спроектирован и внедрён инструмент на базе дистиллированной модели YaGPT~7b:
конвейеры сбора данных и генерации синтетического датасета, дообучение
(Full fine-tune, LoRA, CE~RL), метрики regex и LLM-as-a-judge, инференс-сервис
с оптимизацией FP8 и Eagle. Достигнуты целевые значения F-меры и латентности на контрольных
выборках; эффект на производительность разметки подтверждён внедрением
в платформу разметки ООО «Яндекс.Технологии».%
}

\researchfunding{Работа выполняется на производственной базе ООО «Яндекс.Технологии»
в рамках задачи по улучшению конвейера подготовки данных для обучения
reward-модели системы Alice~AI~LLM~Search.}

\researchpublications{Результаты работы представлены на XV Конгрессе молодых
учёных и планируются к публикации на конференции AINL~2026.}

\maketitle{Бакалавр}

%% Оглавление --- первая страница основного текста (нумерация с~1).
\pagenumbering{arabic}%
\setcounter{page}{1}%
\tableofcontents
\clearpage

%%------------------------------------------------------------
%% ТЕРМИНЫ И ОПРЕДЕЛЕНИЯ
%%------------------------------------------------------------
\chapter*{ТЕРМИНЫ И ОПРЕДЕЛЕНИЯ}\label{chapter:terms}%
\addcontentsline{toc}{chapter}{ТЕРМИНЫ И ОПРЕДЕЛЕНИЯ}%

\noindent\textbf{LLM} (англ. \emph{large language model}, большая языковая модель)~---
генеративная нейросетевая модель, обученная предсказывать следующий токен в тексте
и применяемая для ответов на запросы, суммаризации, классификации и смежных задач.
\textbf{Instruct-модель}~--- языковая модель, дополнительно обученная следовать
инструкциям пользователя и отвечать в заданном формате.

\medskip\noindent\textbf{Информационные контексты}~--- релевантные фрагменты текстов
из внешних источников (веб-документов), на которые опирается поисковая или диалоговая
система при формировании ответа; в работе они служат опорой для проверки фактической
согласованности ответа языковой модели.

\medskip\noindent\textbf{Фактологическая подтверждённость}~--- степень,
в которой утверждение из ответа LLM подтверждается предоставленными
информационными контекстами. В работе эта задача соответствует проверке того,
можно ли считать фрагмент ответа подтверждённым, частично подтверждённым или
неподтверждённым на основании документов.

\medskip\noindent\textbf{Галлюцинация LLM}~--- утверждение, сгенерированное
языковой моделью, но не подтверждаемое исходными контекстами либо противоречащее
им. Для задачи разметки такие ответы опасны, поскольку могут создавать видимость
достоверного факта без опоры на источник.

\medskip\noindent\textbf{Разметка}~--- процесс ручной или полуавтоматической
аннотации данных по заранее заданной схеме; \textbf{исполнитель} (разметчик)~---
специалист, выполняющий разметку в рамках платформы.

\medskip\noindent\textbf{Reward-модель} (модель вознаграждения)~---
обучаемая функция или нейросеть, оценивающая качество ответа (например, степень
подтверждённости фрагмента текста информационными контекстами) числовым сигналом
для обучения с подкреплением или ранжирования вариантов.
\textbf{Reward-функция}~--- отдельное правило или модельная оценка, которая
выдаёт численный сигнал по одному аспекту качества ответа: полноте, атрибуции,
достоверности, фокусу и т.д.

\medskip\noindent\textbf{Дистилляция знаний}~--- перенос поведения крупной
\textbf{модели-учителя} на меньшую по числу параметров \textbf{модель-ученика}
путём обучения ученика на парах «вход~--- ответ учителя» и/или на дополнительных
сигналах распределения вероятностей.

\medskip\noindent\textbf{Дообучение}~--- дополнительное обучение уже предобученной
модели на предметно-специфичных данных. \textbf{Full fine-tune}~--- обновление всех
весов модели; \textbf{LoRA} / \textbf{P-Tuning}~--- параметро-эффективные методы,
при которых основная часть весов заморожена, а адаптация выполняется небольшим
набором добавочных параметров.
\textbf{SFT} (\emph{supervised fine-tuning})~--- дообучение модели на размеченных
парах «вход~--- правильный ответ» с учителем.
\textbf{Дистилляция на уровне ответов} (\emph{response-based distillation})~---
вариант дистилляции, при котором модель-ученик обучается повторять готовые
ответы модели-учителя на заданных входах. \textbf{Модель только с декодером}
(\emph{decoder-only})~--- архитектура трансформера, которая авторегрессионно генерирует текст
последовательно слева направо по циклу.
\textbf{Distillation from large models}~--- дистилляция от крупных моделей:
подход, при котором большая модель-учитель формирует обучающие эталонные ответы или
разметку, а меньшая модель-ученик обучается кросс-энтропийным методом на основе эталонных ответов.

\medskip\noindent\textbf{Синтетический датасет}~--- набор обучающих примеров,
частично или полностью сгенерированных алгоритмами (в т.ч. LLM) по заданным правилам
и ограничениям (искуственно моделируем заданное распределение).
\textbf{Пайплайн}~--- последовательность автоматизированных шагов обработки
данных или модели. \textbf{Чекпоинт модели}~--- сохранённое состояние модели
после этапа обучения, которое можно использовать для инференса (в реальных сервисах) или дальнейшего
дообучения.

\medskip\noindent\textbf{LLM-as-a-judge}~--- метод оценки качества ответов,
при котором другая языковая модель выступает в роли «судьи», сравнивая варианты
или выставляя баллы по формализованным критериям.
\textbf{Поточечная оценка} (\emph{pointwise})~--- оценка одного ответа по шкале критериев без
попарного сравнения с другим ответом.

\medskip\noindent\textbf{CERL} (\emph{Cross Entropy Reinforcement Learning})~---
подход к дообучению политики генерации с использованием кросс-энтропийной функции
потерь и сигналов вознаграждения (в работе~--- стадия выравнивания после SFT).
\textbf{Выравнивание модели} (\emph{alignment})~--- дополнительное обучение,
направленное на согласование поведения модели с заданными критериями качества
и ожиданиями пользователей.

\medskip\noindent\textbf{Инференс}~--- этап применения обученной модели для
получения ответов на входные запросы (в отличие от этапа обучения).
\textbf{Инференс-сервис}~--- программный компонент, предоставляющий доступ к модели,
в работе~--- через \textbf{REST~API}.
\textbf{On-prem}~--- развёртывание сервиса внутри инфраструктуры организации,
без передачи данных во внешние облачные API. \textbf{Production}~--- промышленная
эксплуатационная среда, где сервис доступен пользователям и обслуживает реальные
запросы.

\medskip\noindent\textbf{API} (\emph{application programming interface})~---
программный интерфейс взаимодействия компонентов. \textbf{REST~API}~--- API,
построенный вокруг HTTP-запросов и ресурсов. \textbf{JSON}~--- текстовый формат
структурированных данных, используемый в работе для ответа модели и обмена
данными между сервисами.
\textbf{SLA} (\emph{service level agreement})~--- формализованные требования
к эксплуатационным характеристикам сервиса.
\textbf{Rate limiting}~--- ограничение числа запросов к сервису за единицу
времени.
\textbf{Confidence}~--- численная оценка уверенности модели или сервиса в
сформированном ответе.
\textbf{RPS} (\emph{requests per second})~--- число запросов к сервису в секунду.
\textbf{Retry-After}~--- HTTP-заголовок, сообщающий клиенту, через сколько можно
повторить запрос после какого-либо сбоя.

\medskip\noindent\textbf{Precision, Recall, F-мера}~--- метрики качества классификации
на фиксированном наборе меток: точность (доля верно предсказанных положительных среди
всех предсказанных положительных), полнота (доля найденных положительных среди всех
истинно положительных); \textbf{F-мера}~--- их гармоническое среднее.
\textbf{Базовый вариант}~--- исходное решение или модель, с которыми сравниваются
новые модели или конфигурации. \textbf{Regex-метрики}~--- автоматические проверки
на основе регулярных выражений, которые контролируют наличие нужных полей,
формат ответа и другие структурные признаки.

\medskip\noindent\textbf{Квантизация FP8}~--- представление весов и/или активаций
в формате с плавающей запятой из 8~бит для ускорения вычислений и снижения объёма памяти.
\textbf{BF16}~--- формат чисел с плавающей запятой bfloat16, часто используемый
для обучения и инференса нейросетей. \textbf{GPU}~--- графический процессор,
используемый для ускорения вычислений модели. \textbf{VRAM}~--- видеопамять GPU,
в которой размещаются веса модели, KV-кэш и рабочие буферы.
\textbf{KV-кэш}~--- память ключей и значений механизма внимания трансформера,
которая хранится при генерации и ускоряет продолжение уже начатого ответа.
\textbf{Батч}~--- группа запросов или обучающих примеров, обрабатываемых моделью
за один шаг.

\medskip\noindent\textbf{Спекулятивное декодирование}~--- ускорение генерации за счёт
черновых предсказаний небольшой моделью с последующей верификацией большой моделью;
в работе рассматривается семейство методов \textbf{Eagle}.
\textbf{TTFT} (\emph{time to first token})~--- время от поступления запроса
до генерации первого токена ответа.
\textbf{p95} и \textbf{p99}~--- 95-й и 99-й перцентили задержки: значения,
ниже которых укладываются соответственно 95\% и 99\% запросов.

\medskip\noindent\textbf{Zero-shot}~--- режим применения модели без примеров
решения конкретной задачи в промпте. \textbf{Top-$p$}~--- параметр генерации,
при котором модель выбирает следующий токен из минимального набора наиболее
вероятных токенов с суммарной вероятностью не ниже~$p$.
\textbf{Промпт}~--- текстовая инструкция и входные данные, передаваемые
языковой модели. \textbf{Температура генерации}~--- параметр, управляющий
степенью случайности выбора токенов. \textbf{Top-$k$}~--- ограничение выбора
следующего токена только $k$ наиболее вероятными вариантами.
\textbf{Logits}~--- необработанные числовые оценки вероятностей токенов до
применения нормализации softmax.

\medskip\noindent\textbf{YTSaurus}~--- распределённая система хранения и
обработки данных, используемая для логов, датасетов и промежуточных таблиц.
\textbf{MinHash--LSH}~--- метод приближённого поиска похожих текстов:
MinHash оценивает сходство множеств токенов, а LSH ускоряет поиск почти
дубликатов. \textbf{PII} (\emph{personally identifiable information})~---
персональные данные или сведения, позволяющие идентифицировать человека.
\textbf{Индекс Жаккара}~--- мера сходства двух множеств, равная отношению
размера их пересечения к размеру объединения.

\medskip\noindent\textbf{NLP}~--- обработка естественного языка;
\textbf{CV}~--- компьютерное зрение. \textbf{Crowdsourcing-платформа}~---
платформа, распределяющая задачи между большим числом исполнителей.
\textbf{ML-бэкенд}~--- серверная часть системы, на которой развернута модель машинного.
\textbf{Labeling function}~--- функция, дающая по данным метку.

\medskip\noindent\textbf{HTTP}~--- протокол прикладного уровня для передачи запросов между клиентом
и сервером. \textbf{Stateless-сервис}~--- сервис, не хранящий пользовательское
состояние между запросами.
\textbf{TTL} (\emph{time to live})~--- срок хранения данных или артефакта,
после которого он удаляется автоматически.

\medskip\noindent\textbf{CI/CD}~--- автоматическая сборка, тестирование и выкатка новых версий. \textbf{Docker}~---
технология упаковки приложения в контейнер с зависимостями.
\textbf{Канареечная выкатка}~--- постепенный запуск новой версии на небольшой
доле трафика перед полным развёртыванием.
\textbf{A/B-сравнение}~--- эксперимент, в котором две версии системы сравниваются
на сопоставимых группах запросов или пользователей.
\textbf{wrk} и \textbf{Locust}~--- инструменты нагрузочного тестирования
HTTP-сервисов.

\medskip\noindent\textbf{ZeRO}~--- метод оптимизации памяти при обучении
больших моделей за счёт распределения параметров, градиентов и состояний
оптимизатора. \textbf{Flash-Attention}~--- эффективная реализация механизма
внимания трансформера, уменьшающая потребление памяти и ускоряющая обучение
или инференс.
\textbf{Модель-критик}~--- дополнительная модель в алгоритмах обучения
с подкреплением, оценивающая действия основной модели; в данной работе
отдельная модель-критик не используется.
\textbf{Hugging Face Transformers}~--- библиотека для работы с современными
нейросетевыми языковыми моделями. \textbf{PEFT}~--- семейство параметро-эффективных
методов дообучения, к которому относится LoRA.

\clearpage

%%------------------------------------------------------------
%% ВВЕДЕНИЕ
%%------------------------------------------------------------
\startprefacepage

Современные генеративные языковые модели (LLM) находят широкое применение
в интеллектуальных поисковых системах и голосовых ассистентах. Система
Alice~AI~LLM~Search компании «Яндекс» формирует ответы на поисковые запросы
пользователей, опираясь на \emph{информационные контексты}~--- релевантные
фрагменты текстов, извлечённые из интернет-источников. Качество генерируемых
ответов напрямую определяет пользовательский опыт, поэтому обеспечение их
фактической достоверности является приоритетной задачей~\cite{guo2025deepseek}.
\textbf{Актуальность} темы опирается на обзоры масштабирования и ограничений
LLM~\cite{minaee2024llm-survey,wei2022emergent,touvron2023llama2}, работы по обучению с подкреплением
и выравниванию поведения моделей~\cite{ouyang2022instructgpt,schulman2017ppo,rafailov2023dpo},
методам параметро-эффективного дообучения~\cite{hu2022lora,sanh2019distilbert},
обобщающим исследованиям дистилляции~\cite{hinton2015distilling,distiller-nlp,towards-distillation-theory},
подходам извлечения опоры на документах и оценке согласованности с контекстом~\cite{lewis2020rag,es2023ragas,zhao2023survey},
обзорам инструментов и метрик разметки~\cite{silicon-annotation,automated-annotation-validation,llm-eval-survey},
методу LLM-as-a-judge~\cite{zheng2023judging}, инженерным приёмам ускорения обучения
и инференса~\cite{dao2022flashattention,dao2023flashattention2,rajbhandari2020zero,kwon2023vllm}, квантизации вычислений~\cite{nvidia2022fp8}
и спекулятивному декодированию~\cite{li2024eagle}.

Для обучения \emph{reward-модели}, оценивающей степень подтверждённости
ответа LLM информационными контекстами, требуется большой объём размеченных
данных. Разметчики вручную проверяют каждый фрагмент текста ответа по
трёхклассовой шкале: \textbf{0}~--- не подтверждён, \textbf{1}~---
частично подтверждён, \textbf{2}~--- подтверждён. Текущий процесс занимает
в среднем 7–8~минут на задание, из которых 3–4~минуты на задание расходуется
на поиск релевантных документов в информационных контекстах и ещё 4~минуты на задание~---
на верификацию фактов. Подобная трудоёмкость делает разметку «бутылочным
горлышком» в конвейере подготовки данных~\cite{silicon-annotation}.

Для частичной автоматизации поиска контекстов текущий пайплайн использует
модель DeepSeek-V3 в качестве внешнего LLM-ассистента. Однако этот подход
имеет существенные недостатки: время генерации ответа составляет 2–5~минут
на 1~фрагмент из задания,
применяются неструктурированные промпты, а качество ответов удовлетворяет
исполнителей лишь в~50\% случаев. Ключевые метрики текущей разметки
(F-мера~$\approx 0{,}60$, Precision~$\approx 0{,}55$, Recall~$\approx 0{,}57$)
не отвечают требованиям к обучающим данным для production-моделей~\cite{automated-annotation-validation}.

Для устранения указанных проблем предлагается разработать
\textbf{специализированный инструмент} на основе дистиллированной языковой
модели~\cite{hinton2015distilling}. Дистиллированная модель меньших размеров
(7–32~млрд параметров вместо $\sim600$~млрд у DeepSeek-V3), обученная
непосредственно на задаче анализа информационных контекстов, обеспечит
значительно меньшее время инференса и более высокое качество ответов
благодаря узкоспециализированному дообучению.

\textbf{Цель работы}~--- улучшить систему разметки фрагментов, ускорив процесс
\textbf{в два раза} до~3--4~минут на задание (сократив время поиска контекстов
вдвое), повысив точность до~F-меры $> 0{,}85$ (с текущих~$\approx 0{,}60$),
обеспечив быстрый отклик LLM с латентностью~$\le 30$~секунд на 1~фрагмент
из задания (вместо~2--5~минут на 1~фрагмент из задания)
при помощи \textbf{спекулятивного декодирования} (алгоритм Eagle) и
\textbf{квантизации в~FP8}, минимизировав галлюцинации (с~$\approx 30$\%) и
подняв удовлетворённость команды выше~70\% (с~$\approx 50$\%), сделав её
быстрой, точной, надёжной и удобной.

\textbf{Задачи:}
\begin{enumerate}
    \item Подобрать архитектуру и размер генеративной модели-«ученика»
          с балансом скорость/качество и спроектировать процедуру её обучения.
    \item Сформировать синтетические датасеты запросов и ответов
          (distillation from large models), реализовать очистку/фильтрацию
          и контроль распределений.
    \item Реализовать пайплайны оркестрации (статический и динамический граф)
          для подготовки данных, генерации, обучения, оценки и инференса.
    \item Разработать метрики автоматической оценки качества ответов
          (правила/регулярные выражения + LLM-as-a-judge) и провести серию
          экспериментов по гиперпараметрам и оптимизациям обучения/инференса.
\end{enumerate}

\textbf{Объект исследования}~--- процесс текстовой разметки ответов
генеративной языковой модели Alice~AI~LLM~Search по критерию достоверности
фактов.

\textbf{Предмет исследования}~--- методы и алгоритмы автоматического анализа
информационных контекстов для ускорения и повышения качества текстовой
разметки.

\textbf{Новизна работы.} В отличие от существующих решений (универсальные
LLM-API, платформы разметки, DeepSeek-V3 без дообучения), предлагаемый
подход сочетает три элемента, ранее не применявшихся совместно в задаче
ассистирования разметчиков:
\begin{enumerate}
    \item \emph{Специализированная дистилляция} ответов учителя
    (DeepSeek-V3.1-Terminus) на корпусе реальных доменных запросов
    исполнителей с балансировкой по типам запросов;
    \item \emph{Стадия выравнивания методом CE~RL} с шестью факторными
    функциями вознаграждения (атрибуция, полнота, достоверность,
    ясность, цитаты, фокус), оцениваемыми LLM-судьёй, и
    их групповой агрегацией;
    \item \emph{On-prem инференс} специализированной модели (YaGPT~7b)
    со структурированным выходом, встроенной в интерфейс платформы
    разметки — без передачи данных во внешние API.
\end{enumerate}
Совокупность этих элементов позволяет впервые достичь латентности
$\le 30$~с на 1~фрагмент из задания при F-мере $> 0{,}85$ в задаче разметки фактологической
подтверждённости в закрытой производственной инфраструктуре.

\textbf{Практическая значимость.} Инструмент внедрён в промышленную платформу
текстовой разметки ООО «Яндекс.Технологии» в виде диалогового окна и доступен
через REST~API. Ожидаемый измеримый эффект: ускорение разметки в~2~раза (до~3--4~мин на задание),
F-мера выше~$0{,}85$, латентность ответа модели~$\le 30$~с на 1~фрагмент из задания, рост удовлетворённости
исполнителей выше~70\%, снижение доли несогласованных с контекстом ответов
по сравнению с универсальной внешней LLM.
Внедрение подтверждается актом о внедрении от ООО «Яндекс.Технологии».
Результаты работы представлены на XV Конгрессе молодых учёных; материалы
планируются к публикации на конференции AINL~2026.

\textbf{Структура работы.} В первой главе приведён аналитический обзор методов
ускоренной разметки и оценки достоверности ответов LLM: рассмотрены
существующие платформы разметки, подходы на основе LLM-ассистентов, методы
дистилляции и дообучения, метрики качества и выбор технологического стека;
сформулированы требования и постановка задачи. Во второй главе описано
проектирование инструмента анализа информационных контекстов~--- общая
архитектура (C4), конвейеры сбора данных, генерации датасета и обучения,
архитектура инференс-сервиса, логическая модель данных. В третьей главе
приведены реализация, тестирование и оптимизация инструмента: конвейеры,
синтетический датасет, эксперименты по дообучению и стадия CE~RL,
регрессионные и нагрузочные проверки, результаты внедрения, а также
инженерная оптимизация инференс-сервиса (квантизация FP8, спекулятивное
декодирование Eagle), методика замеров и сравнение конфигураций.

%%------------------------------------------------------------
%% ГЛАВА 1 — АНАЛИТИЧЕСКИЙ ОБЗОР МЕТОДОВ УСКОРЕННОЙ РАЗМЕТКИ
%%------------------------------------------------------------
\chapter{Аналитический обзор методов ускоренной разметки и оценки достоверности ответов LLM}

\startrelatedwork

\section{Задача оценки достоверности ответов LLM}

С ростом применения LLM в поисковых системах остро встала проблема
\emph{галлюцинаций}~--- генерации фактически неверной информации~\cite{minaee2024llm-survey,wei2022emergent}.
Для контроля достоверности ответов системы Alice~AI~LLM~Search
применяется отдельная reward-модель, оценивающая степень
подтверждённости каждого фрагмента ответа информационными
контекстами~\cite{zhao2023survey,lewis2020rag}. Обучение такой модели требует
сотен тысяч размеченных образцов.

Разметка каждого образца~--- это задача \emph{многоклассовой классификации}:
разметчик изучает фрагмент текста ответа LLM и список информационных
контекстов, а затем присваивает класс:
\textbf{0}~--- факт не подтверждён ни одним контекстом,
\textbf{1}~--- факт частично подтверждён ($>50\%$ сходимости),
\textbf{2}~--- факт полностью подтверждён.
Ключевая трудоёмкость операции~--- поиск среди десятков документов именно тех
фрагментов, которые подтверждают или опровергают конкретный факт.

\section{Анализ существующих инструментов и подходов}

\subsection{Платформы разметки общего назначения}

На рынке существует ряд зрелых платформ разметки данных: Label Studio,
Argilla, Prodigy, Scale~AI Nucleus, Toloka (Яндекс). Все они предоставляют
пользовательский интерфейс разметчика и инструменты контроля качества,
однако принципиально не решают проблему поиска релевантных фрагментов
в контекстах. Разметчик по-прежнему просматривает документы вручную.

\textbf{Label Studio}~\cite{labelstudio} — открытая платформа с поддержкой
широкого спектра типов разметки (NLP, CV, audio). Позволяет добавить
ML-бэкенд для пре-аннотации, однако не содержит специализированных
компонентов для задачи оценки фактологической подтверждённости: поиска фрагментов документа,
подтверждающих или опровергающих конкретное утверждение.

\textbf{Argilla}~\cite{argilla} — платформа, ориентированная на сбор
обратной связи для RLHF. Поддерживает обратную связь на уровне ответов
модели, но не даёт разметчику инструмента для грануларной работы
с контекстами: нет механизма выделить конкретный фрагмент-источник
и сопоставить его с утверждением в ответе LLM.

\textbf{Snorkel AI}~\cite{snorkel} предлагает программируемую разметку
(labeling functions) и слабый надзор. Этот подход эффективен для задач
с чёткими паттернами, однако верификация фактических утверждений
в ответах LLM требует семантического понимания, которое не формализуется
детерминированными правилами.

\textbf{Scale AI / Toloka} — crowdsourcing-платформы. Масштабируют ручной
труд, но не сокращают время на одного исполнителя и сохраняют требование
к предметной экспертизе.

\textbf{Вывод}: ни одна из платформ не предоставляет специализированного
ассистента, который бы \emph{автоматически} находил и атрибутировал
подтверждающие фрагменты из контекстов. Расширение любой из них до
требуемой функциональности потребовало бы разработки того же
ML-компонента, что и в данной работе, плюс дополнительной интеграции
в закрытую инфраструктуру Яндекса.

\subsection{Подходы на основе LLM-ассистентов}

Второй класс решений~--- использование общей инструктируемой LLM
в качестве ассистента разметчика. В таблице~\ref{tab:llm-comparison}
приведено сравнение по ключевым параметрам задачи.

\begin{table}[h]
\caption{Сравнение LLM-ассистентов для разметки фактологической подтверждённости}\label{tab:llm-comparison}
\centering
\small
\begin{tabularx}{\textwidth}{|l|>{\centering}X|>{\centering\arraybackslash}X|}
\hline
\textbf{Решение} & \textbf{Латент-ность} & \textbf{F-мера} \\
\hline
GPT-4o (API) & 5–15~с & $\sim0{,}72$\textsuperscript{*} \\
\hline
Claude 3.5 Sonnet (API) & 8–20~с & $\sim0{,}70$\textsuperscript{*} \\
\hline
DeepSeek-V3 (текущий) & 2–5~мин на 1~фрагмент из задания & $0{,}60$ \\
\hline
YaGPT~7b zero-shot & 20–40~с & $0{,}51$ \\
\hline
\textbf{YaGPT~7b + SFT + CE~RL} & \textbf{$\le$30~с на 1~фрагмент из задания} & \textbf{$>$0{,}85} \\
\hline
\multicolumn{3}{l}{\footnotesize * Оценка по аналогичным задачам fact-checking из~\cite{silicon-annotation, llm-eval-survey}; прямое тестирование} \\
\multicolumn{3}{l}{\footnotesize на закрытых данных не проводилось.} \\
\end{tabularx}
\end{table}

Публичные API (GPT-4o, Claude) недоступны как on-prem решения: данные
разметчиков содержат внутреннюю информацию компании. DeepSeek-V3 в
текущей конфигурации работает без специализированного промпта и
дообучения~--- отсюда 2–5~минут на 1~фрагмент из задания и F-мера~$= 0{,}60$.
YaGPT~7b в режиме zero-shot даёт приемлемую латентность, но не умеет
возвращать структурированный ответ с атрибуцией и цитированием:
F-мера~$= 0{,}51$.

\textbf{Ключевой вывод}: разрыв между zero-shot YaGPT~7b ($0{,}51$)
и целевым уровнем ($>0{,}85$) не закрывается промпт-инжинирингом~---
модель не видела доменных примеров и не умеет формировать структурированные
ответы с атрибуцией. Дообучение на задаче является необходимым и
единственным on-prem-совместимым путём.

\subsection{Анализ задачи и формулировка требований}

На основании сравнения определены ключевые требования к разрабатываемому
инструменту:
\begin{enumerate}
    \item \textbf{Латентность~$\le 30$~с на 1~фрагмент из задания}
          (текущий узкий участок~--- 2–5~мин на 1~фрагмент из задания);
    \item \textbf{F-мера~$> 0{,}85$} при трёхклассовой классификации достоверности;
    \item \textbf{On-prem деплой}: модель работает на GPU Яндекса без
          передачи данных во внешние API;
    \item \textbf{Структурированный ответ} с атрибуцией к конкретным
          фрагментам контекстов (attribution, quotes);
    \item \textbf{Специализация}: дообучение на запросах исполнителей
          платформы разметки, а не на общем корпусе.
\end{enumerate}

Ни одно из рассмотренных решений не удовлетворяет всем пяти требованиям
одновременно. Данная работа закрывает этот пробел.

\subsection{Сравнение подходов к автоматизации разметки (итог)}

В таблице~\ref{tab:approaches} сведено итоговое сравнение по классам
подходов.

\begin{table}[h]
\caption{Сравнение классов подходов к автоматизации разметки}\label{tab:approaches}
\centering
\begin{tabularx}{\textwidth}{|l|>{\centering}X|>{\centering\arraybackslash}X|}
\hline
\textbf{Подход} & \textbf{Скорость} & \textbf{Качество} \\
\hline
Ручная разметка & 7–8~мин на задание & Высокое \\
\hline
Платформы (Label Studio и др.) & 7–8~мин на задание & Высокое \\
\hline
Универсальный LLM-API & 5–20~с & Среднее \\
\hline
DeepSeek-V3 без FT (текущий) & 2–5~мин на 1~фрагмент из задания & Низкое \\
\hline
\textbf{YaGPT~7b + SFT + CE~RL} & \textbf{$\le$30~с на 1~фрагмент из задания} & \textbf{Высокое} \\
\hline
\end{tabularx}
\end{table}

\section{Методы дистилляции языковых моделей}

Дистилляция знаний~\cite{hinton2015distilling}~--- метод переноса знаний
из большой модели-учителя в меньшую модель-ученика. В классическом
варианте ученик обучается воспроизводить «мягкие» вероятностные
распределения учителя, а не жёсткие метки классов, что позволяет передать
дополнительную информацию о межклассовых отношениях.

Для языковых моделей распространена \textbf{дистилляция на уровне ответов}
(\emph{response-based distillation}): ученик обучается на корпусе
вопрос-ответных пар, где ответы сгенерированы учителем~\cite{sanh2019distilbert}.
Этот подход особенно эффективен, когда имеется доступ к сильному учителю,
но не к его параметрам (black-box distillation)~\cite{distiller-nlp}.

В данной работе применяется метод \textbf{Distillation from Large Models}:
модель DeepSeek-V3.1-Terminus генерирует эталонные ответы на запросы
исполнителей, которые затем используются для дообучения модели-ученика
(YaGPT 7b/8b/32b). Основные этапы: сбор промптов $\to$ генерация ответов
учителем $\to$ пост-обработка (фильтрация, балансировка) $\to$ обучение ученика.

\section{Методы дообучения языковых моделей}

\textbf{Full fine-tuning} предполагает обновление всех параметров модели.
Позволяет добиться максимальной специализации, однако требует значительного
объёма GPU-памяти и большого датасета ($\ge 10{,}000$ образцов).
Используется в начальных экспериментах на синтетическом датасете.

\textbf{LoRA} (Low-Rank Adaptation)~\cite{hu2022lora}~--- эффективный метод
параметрического дообучения: добавляются низкоранговые адаптерные матрицы
к весам трансформера, обучаются только они (0{,}1–1\% параметров).
Применяется для дообучения на реальных запросах исполнителей ($<1{,}500$ образцов).

\textbf{RLHF / CERL} (Cross Entropy Reinforcement Learning)~\cite{ouyang2022instructgpt, schulman2017ppo}~---
метод обучения с подкреплением, где reward-сигнал поступает от LLM-судьи.
Планируется применить после получения базового чекпоинта для улучшения
качества структурированности ответов.

\textbf{DPO} (Direct Preference Optimization)~\cite{rafailov2023dpo}~--- альтернатива
RLHF без явного обучения reward-модели. Рассматривается как запасная
стратегия при нестабильности CERL.

\section{Методы оценки качества языковых моделей}

Традиционные метрики (BLEU, ROUGE) плохо коррелируют с человеческой
оценкой для задач генерации структурированных текстов~\cite{llm-eval-survey}.
В работе используется комбинация двух подходов.

\textbf{Regex-метрики}~--- проверка соответствия ответа модели-ученика
заданному формату с помощью регулярных выражений. Оценивают корректность
структуры ответа: наличие обязательных полей, формат ссылок на источники,
корректность классификационных меток. Работают быстро, детерминированы.

\textbf{LLM-as-a-judge}~\cite{zheng2023judging}~--- оценка качества ответа
с помощью более мощной модели (DeepSeek-V3.1-Terminus). Промпт содержит
ответ учителя, ответ ученика и критерии оценки: точность фактов,
глубина анализа, наличие галлюцинаций, соответствие контексту.
Позволяет оценить семантическое качество ответов.

\section{Выбор технологического стека}

Выбор технологий обусловлен производственными ограничениями
ООО «Яндекс.Технологии» и требованиями к производительности.

\textbf{PyTorch}~--- основной фреймворк машинного обучения. Выбран за развитую
экосистему (Hugging Face Transformers, PEFT), поддержку
Flash-Attention~\cite{dao2022flashattention, dao2023flashattention2}
и ZeRO-оптимизации~\cite{rajbhandari2020zero} для эффективного обучения
на GPU Nvidia H100.

\textbf{Внутренний аналог Apache Airflow (статический)}~--- платформа
оркестрации конвейеров с фиксированным графом задач. Используется для
сбора данных, генерации датасета и обучения модели.

\textbf{Внутренний аналог Apache Airflow (динамический)}~--- платформа
исполнения конвейеров в реальном времени с динамическим графом узлов.
Используется для обработки запросов в production.

\textbf{YTSaurus}~--- платформа распределённого хранения и обработки данных
с открытым исходным кодом. Используется для хранения логов, датасетов
и артефактов обучения.

\textbf{Внутренний аналог vLLM}~\cite{kwon2023vllm}~--- фреймворк ускоренного
инференса с~PagedAttention. Используется для деплоя обученной модели
в production.

\textbf{NLTK}~--- используется для предобработки текстовых данных:
токенизация, стемминг, анализ распределения токенов.

\finishrelatedwork

\section{Защита данных, ограничения логов и риски применения LLM}

Конвейер обучения и инференса опирается на логи запросов исполнителей
и фрагменты ответов LLM. Для снижения юридических и репутационных рисков
приняты следующие принципы. \textbf{Минимизация данных:} в таблицы
YTSaurus попадают только поля, необходимые для воспроизведения промпта
и оценки качества (идентификатор задания, тип запроса, нормализованный
текст, ссылки на контексты без полного HTML-страниц, где это возможно).
\textbf{Псевдонимизация:} идентификаторы исполнителей заменяются на
устойчивые хэши; прямые ФИО и контактные данные в обучающих выборках
не хранятся. \textbf{Сроки хранения:} сырые логи и промежуточные
артефакты пайплайна удаляются по политике TTL оркестратора; чекпоинты
и финальные датасеты версионируются и доступны ограниченному кругу
сервисных аккаунтов.

\textbf{Риски использования LLM} в цикле разметки включают: утечку
чувствительных формулировок из контекстов в логи ошибок инференса;
«запоминание» редких шаблонов из обучающей выборки (меморизация);
смещение распределения ответов при дообучении на узком корпусе.
Противодействие: on-prem деплой без внешних API, фильтрация PII
на этапе предобработки, регулярный пересмотр выборки и контроль
метрик на отложенной выборке, не участвовавшей в подборе гиперпараметров.

\chapterconclusion

В результате аналитического обзора установлено, что основная проблема
действующего процесса разметки связана не только с собственно принятием
решения по метке достоверности, но и с длительным поиском подтверждающих
контекстов, проверкой противоречий и ожиданием ответа универсальной внешней
LLM. Существующие платформы разметки общего назначения (Label Studio,
Argilla, Snorkel~AI) хорошо закрывают задачи организации рабочего процесса,
но не решают доменную задачу оценки фактологической подтверждённости
для фрагментов ответов LLM без
дополнительного интеллектуального ассистента. Подход с использованием
крупной внешней модели без дообучения также недостаточен: он даёт высокую
латентность, не гарантирует стабильный структурированный формат ответа и
не обеспечивает требуемую F-меру выше~$0{,}85$.

Анализ методов показал, что наиболее обоснованным направлением является
разработка специализированной дистиллированной модели: крупная модель-учитель
формирует эталонные ответы, а компактная модель-ученик обучается на целевом
домене и затем разворачивается on-prem внутри производственного контура.
Для достижения требуемого качества целесообразно сочетать несколько этапов:
Full fine-tune на синтетическом корпусе, параметро-эффективное дообучение
LoRA на реальных запросах и последующую стадию CE~RL для выравнивания
структуры, полноты, атрибуции и фокуса ответов. Такая схема лучше соответствует
ограничениям задачи, чем только настройка промптов или использование
универсального API.

В качестве методов контроля качества выбрана комбинация детерминированных
regex-проверок и оценки LLM-as-a-judge, поскольку одна группа метрик
проверяет структурированный ответ генеративной модели, заданный в чётком
формате при помощи системного промпта, а другая позволяет оценить семантические
свойства: достоверность, полноту, атрибуцию, качество цитирования и фокус.
Выбранный технологический стек (PyTorch, YTSaurus, внутренний аналог
vLLM, оркестраторы статического и динамического графа) согласуется с
производственными требованиями к воспроизводимости, хранению данных,
скорости инференса и изоляции чувствительных контекстов. Тем самым в первой
главе обоснована постановка задачи, целевые метрики и необходимость
проектирования специализированной системы, а не отдельного экспериментального
скрипта или разовой интеграции с внешней LLM.

%%------------------------------------------------------------
%% ГЛАВА 2 — ПРОЕКТИРОВАНИЕ ИНСТРУМЕНТА АНАЛИЗА КОНТЕКСТОВ
%%------------------------------------------------------------
\chapter{Проектирование инструмента анализа информационных контекстов}

\section{Общая архитектура системы}

На рисунке~\ref{fig:containers} приведена контейнерная диаграмма системы
(C4 Level~2~--- Container): границы Annotation Assistant, основные контейнеры
и потоки данных относительно исполнителя и платформы разметки; красным выделен
вклад автора.
Инструмент анализа информационных контекстов выступает промежуточным
звеном между исполнителем и платформой разметки: принимает запрос
исполнителя и информационные контексты, возвращает структурированный
анализ достоверности.

\medskip\noindent\textbf{Диаграммы в нотациях C4/UML и логическая модель данных.}\par
Для проектной документации используются диаграммы \textbf{C4 Container}
(рис.~\ref{fig:containers}),
\textbf{UML Sequence} для сценария вызова ассистента (рис.~\ref{fig:sequence}),
а также блок-схемы конвейеров на рисунках~\ref{fig:data-pipeline}--\ref{fig:train-pipeline},
отражающие потоки данных между хранилищами и сервисами и выполняющие роль
упрощённых \textbf{DFD}. Полное BPMN-описание допускается как расширение
тех же потоков при согласовании с заказчиком; в тексте ВКР достаточно
показанной декомпозиции этапов.

На концептуальном уровне (\textbf{ER}) выделяются сущности:
\emph{Исполнитель}, \emph{Задание разметки}, \emph{Информационный контекст},
\emph{Запрос к ассистенту}, \emph{Структурированный ответ модели},
\emph{Запись обучающего датасета}, \emph{Версия чекпоинта}.
Исполнитель выполняет множество заданий; задание объединяет множество контекстов;
запрос к ассистенту порождает один ответ с атрибуцией; записи датасета
агрегируются из логов и ответов учителя; чекпоинт версионируется и связывается
с конфигурацией сервиса. Физическое отображение~--- таблицы и каталоги
YTSaurus, объектное хранилище артефактов (см.~рис.~\ref{fig:containers},
приложение~\ref{app:diagrams}).

\begin{figure}[h]
\caption{Контейнерная диаграмма системы (C4 Level 2 — Container). Красным выделен вклад автора}\label{fig:containers}
\centering
\includegraphics[width=\textwidth]{diagrams/best_diagram_ever.png}
\end{figure}

На рисунке~\ref{fig:sequence} представлена диаграмма последовательности
запроса анализа информационных контекстов: от действия исполнителя до
получения структурированного ответа (класс достоверности, атрибуция, цитаты).

\begin{figure}[h]
\caption{Диаграмма последовательности: запрос анализа информационных контекстов}\label{fig:sequence}
\centering
\includegraphics[width=\textwidth]{diagrams/Sequence_AnalyzeRequest.png}
\end{figure}

Компонентная диаграмма конвейера обучения модели приведена
в Приложении~\ref{app:diagrams}.

\section{Конвейер сбора и предобработки данных}\label{sec:data-pipeline}

Конвейер реализован в виде последовательной цепочки узлов
на платформе статического оркестратора (рисунок~\ref{fig:data-pipeline}).
Входными данными служат логи платформы разметки, хранящиеся в~YTSaurus.

\begin{figure}[h]
\caption{Конвейер сбора и предобработки данных}\label{fig:data-pipeline}
\centering
\begin{tikzpicture}[
    pnode/.style={draw, rectangle, rounded corners=4pt, fill=orange!15,
                 minimum width=2.4cm, minimum height=0.9cm,
                 text width=2.2cm, align=center, font=\small},
    arr/.style={-{Stealth}, thick},
    node distance=0.5cm
]
\node[pnode] (logs)   at (0,0)     {Логи\\YTSaurus};
\node[pnode] (clean)  at (2.9,0)   {Очистка\\и норм.};
\node[pnode] (dedup)  at (5.8,0)   {Дедупли-\\кация};
\node[pnode] (types)  at (8.7,0)   {Анализ\\типов};
\node[pnode] (out)    at (11.6,0)  {Датасет\\запросов};

\draw[arr] (logs)  -- (clean);
\draw[arr] (clean) -- (dedup);
\draw[arr] (dedup) -- (types);
\draw[arr] (types) -- (out);
\end{tikzpicture}
\end{figure}

Узел \textbf{«Очистка и нормализация»}: удаление дублирующих пробелов,
html-тегов, нормализация unicode, отфильтровка слишком коротких
запросов ($<10$ токенов) и нетематических сообщений.

Узел \textbf{«Дедупликация»}: удаление точных и нечётких дубликатов
(с порогом Jaccard~$\ge 0{,}9$) с помощью MinHash LSH.

Узел \textbf{«Анализ типов запросов»}: кластеризация запросов (K-means
на TF-IDF эмбеддингах) для последующего контроля равномерности
распределения в синтетическом датасете.

\section{Конвейер генерации синтетического датасета}\label{sec:synth-design}

Конвейер генерации (рисунок~\ref{fig:synth-pipeline}) состоит из двух
логических ветвей: синтез доменных запросов по типологиям реальных логов
и получение эталонных ответов моделью-учителем с последующей фильтрацией.
Ниже зафиксированы принятые проектные решения; реализация и численные
итоги приведены в главе~3.

\begin{figure}[h]
\caption{Конвейер генерации синтетического датасета}\label{fig:synth-pipeline}
\centering
\begin{tikzpicture}[
    pnode/.style={draw, rectangle, rounded corners=4pt, fill=green!12,
                 minimum width=2.6cm, minimum height=1cm,
                 text width=2.4cm, align=center, font=\small},
    arr/.style={-{Stealth}, thick},
    node distance=0.5cm
]
\node[pnode] (types)  at (0,0)     {Типы\\запросов};
\node[pnode] (prompt_q) at (3.0,0) {Промпт для\\генерации\\запросов};
\node[pnode] (gen_q)  at (6.0,0)   {DeepSeek\\генерирует\\запросы};
\node[pnode] (filter) at (9.0,0)   {Фильтрация\\(распред.)};
\node[pnode] (prompt_a) at (0,-2.5){Промпт для\\генерации\\ответов};
\node[pnode] (gen_a)  at (3.0,-2.5){DeepSeek\\генерирует\\ответы};
\node[pnode] (postproc) at (6.0,-2.5){Пост-\\обработка};
\node[pnode] (dataset)at (9.0,-2.5){Датасет\\($\ge$10k)};

\draw[arr] (types)   -- (prompt_q);
\draw[arr] (prompt_q) -- (gen_q);
\draw[arr] (gen_q)   -- (filter);
\draw[arr] (filter.south) -- ++(0,-0.5) -| (prompt_a.north);
\draw[arr] (prompt_a) -- (gen_a);
\draw[arr] (gen_a)   -- (postproc);
\draw[arr] (postproc) -- (dataset);
\end{tikzpicture}
\end{figure}

\subsection{Назначение синтетического корпуса и риски}

Пары «запрос исполнителя~--- эталонный разбор контекстов» необходимы,
чтобы в supervised-фазе задать объём и стиль поведения модели-ученика
в домене платформы разметки при ограниченном числе вручную размеченных
экземпляров~\cite{silicon-annotation,automated-annotation-validation}.
Согласно парадигме дистилляции на уровне ответов учитель (DeepSeek-V3.1-Terminus)
задаёт целевые ответы на масштабе, недоступном ручной разметке~\cite{hinton2015distilling}.
Риск смещения распределения при чисто синтетическом обучении учитывается:
после SFT на синтетике выполняется дообучение LoRA на реальных логах и
стадия выравнивания, что согласуется с рекомендациями по оценке и смешению
источников данных для LLM~\cite{llm-eval-survey}.

\subsection{Типология, кластеры и квотирование}

Опорными остаются четыре типа запросов, выделенные по логам
(см.~перечисление в главе~3). На этапе проектирования зафиксировано:
\begin{itemize}
    \item для каждого типа в few-shot шаблон включаются \textbf{3--5}
          репрезентативных примеров из соответствующих кластеров K-means
          на TF-IDF (согласованно с узлом «Анализ типов» конвейера сбора);
    \item целевые доли типов в финальном корпусе задаются по эмпирическим
          частотам после дедупликации логов;
    \item генерация организуется \textbf{стратифицированно}: на каждой
          итерации пересчитывается текущий дефицит типа относительно квоты,
          и объём запросов следующей волны распределяется пропорционально
          дефициту.
\end{itemize}
Отбраковываются дословные и почти дословные повторы few-shot примеров,
а также кандидаты ниже порога по числу токенов (согласованного с
порогом $\ge 10$ токенов на этапе очистки логов).

\subsection{Промптинг генерации запросов}

Инструкция к модели-генератору запросов формулируется как задача
«синтезировать новые формулировки той же семантической цели», без
копирования лексики примеров; допускается вариация длины и наличие
уточняющих подвопросов внутри типа. Параметры декодирования на стороне
генератора запросов (температура, top-$p$, лимит выхода) подбираются
эмпирически на отложенной мини-выборке: цель~--- разнообразие формулировок
при сохранении доменной принадлежности к задаче оценки фактологической
подтверждённости.

\subsection{Промптинг и постобработка ответов учителя}

Промпт к DeepSeek-V3.1-Terminus строится в формате, близком к
production-разметке: \textbf{zero-shot} относительно конкретной пары
«запрос~--- контексты», но с полностью зафиксированной схемой выхода.
Целевая структура включает поля \texttt{relevant\_contexts},
\texttt{analysis}, \texttt{conclusion}, \texttt{label}
(в промышленном структурированном формате, эквивалентном листингу~\ref{lst:gen-prompt}).
Постобработка проектируется как конвейер проверок:
\begin{itemize}
    \item разбор структурированного ответа и отсечение ответов вне схемы;
    \item проверка индексов контекстов на принадлежность диапазону;
    \item regex- и эвристические фильтры против пустого анализа,
          несоответствия числу ссылок и заведомо неконсистентных меток;
    \item контроль баланса классов $0/1/2$ на уровне батча (при перекосе
          следующие вызовы учителя приоритизируют недостающий класс).
\end{itemize}

\subsection{Статистический контроль и версионирование}

Согласование наблюдаемых долей типов запросов в синтетическом корпусе
с долями в предобработанных логах проверяется $\chi^2$-критерием на
уровне значимости $0{,}05$; при отвержении гипотезы планируется
добор генерации по недопредставленным типам либо корректировка квот.
Каждая выпущенная версия корпуса сопровождается в метаданных YTSaurus
идентификатором запуска пайплайна, версией промптов и параметров
декодирования учителя, чтобы обеспечить воспроизводимость экспериментов
по обучению и аудит регрессий качества.

\section{Конвейер обучения модели}

На рисунке~\ref{fig:train-pipeline} представлена схема конвейера обучения,
реализованного на статическом оркестраторе.

\begin{figure}[h]
\caption{Конвейер обучения модели}\label{fig:train-pipeline}
\centering
\begin{tikzpicture}[
    pnode/.style={draw, rectangle, rounded corners=4pt, fill=blue!10,
                 minimum width=3cm, minimum height=1cm,
                 text width=2.8cm, align=center, font=\small},
    arr/.style={-{Stealth}, thick},
    node distance=0.4cm
]
\node[pnode] (data)    at (0,0)    {Датасет\\(запросы +\\ответы учителя)};
\node[pnode] (ft)      at (3.8,0)  {Full fine-tune\\(синт. датасет)};
\node[pnode] (lora)    at (7.6,0)  {LoRA/P-Tuning\\(реальные данные)};
\node[pnode] (cerl)    at (11.4,0) {CERL\\(LLM-судья)};

\node[pnode] (eval)    at (5.7,-2.5){Оценка:\\regex +\\LLM-judge};
\node[pnode] (best)    at (9.5,-2.5){Лучший\\чекпоинт};

\draw[arr] (data) -- (ft);
\draw[arr] (ft)   -- (lora);
\draw[arr] (lora) -- (cerl);
\draw[arr] (ft.south)  -- ++(0,-0.6) -- (eval.north west);
\draw[arr] (lora.south) -- ++(0,-0.6) -- (eval.north);
\draw[arr] (cerl.south) -- ++(0,-0.6) -- (eval.north east);
\draw[arr] (eval) -- (best);
\end{tikzpicture}
\end{figure}

Сетка гиперпараметров:
\begin{itemize}
    \item Число эпох: 3, 5, 10;
    \item Learning rate: $10^{-6}$, $10^{-5}$, $10^{-4}$;
    \item Стратегии scheduler: \texttt{cosine\_annealing}, \texttt{reduce\_on\_plateau};
    \item Оптимизации инфраструктуры: Flash-Attention~2/3~\cite{dao2023flashattention2},
          ZeRO-2/3~\cite{rajbhandari2020zero}, gradient checkpointing;
    \item GPU: Nvidia H100~80GB, количество определяется бенчмарком.
\end{itemize}

\section{Архитектура инференс-сервиса}

Сервис развёрнут на основе внутреннего аналога vLLM~\cite{kwon2023vllm}
с PagedAttention для эффективного управления KV-кэшем.
Запрос исполнителя поступает в динамический пайплайн оркестратора,
где составляется промпт (информационные контексты + вопрос исполнителя)
и отправляется в инференс-сервер. Ответ возвращается в структурированном
формате, заданном системным промптом, и отображается в диалоговом окне
платформы разметки.

REST API сервиса:
\begin{itemize}
    \item \texttt{POST /api/v1/analyze} — основной эндпоинт;
    \item Тело запроса: \texttt{\{query: str, contexts: list[str]\}};
    \item Тело ответа: \texttt{\{label: int, analysis: str,
          relevant\_contexts: list[int], confidence: float\}};
    \item Целевое время ответа: $\le 30$~сек на 1~фрагмент из задания.
\end{itemize}

\section{Нефункциональные требования и эксплуатация сервиса}

Помимо функциональной спецификации API зафиксированы \textbf{SLA}
для production: доступность сервиса анализа не ниже~99{,}5\% в месяц
(исключая плановые окна обновления модели не более двух часов
в квартал); p95 задержки ответа при типичной длине контекста
(до~32\,000 токенов в сумме по полям) не превышает~35~сек;
p99~--- не более~45~сек. При превышении порогов включается эскалация
в дежурную смену платформы разметки.

\textbf{Пропускная способность и очереди.} API защищён
rate limiting по ключу платформы и по пользователю: не более
$N$ параллельных запросов на пул GPU, очередь отклоняется с кодом
HTTP~429 и заголовком \texttt{Retry-After}. Для длинных запросов
поддерживается асинхронный режим (идентификатор задачи + опрос статуса),
чтобы не блокировать воркеры веб-интерфейса.

\textbf{Наблюдаемость.} Каждый запрос получает \texttt{trace\_id};
в структурированные логи пишутся длительность этапов пайплайна
(сборка промпта, вызов инференса, постобработка структурированного ответа),
код ответа
модели, число принятых токенов при спекулятивном декодировании,
размер батча на GPU. Метрики экспортируются в систему мониторинга
(графики задержки, ошибок парсинга, доли отказов черновика Eagle).
Алерты на рост доли HTTP~5xx и на падение доли успешных regex-проверок
структуры ответа.

\textbf{Отказоустойчивость.} Развёрнуто не менее двух реплик
stateless-слоя API за балансировщиком; при недоступности одной ноды
инференса трафик переключается на резервную. Чекпоинт модели хранится
в объектном хранилище с версионированием; откат на предыдущую версию
выполняется без пересборки образа контейнера.

\chapterconclusion

Во второй главе разработана архитектура инструмента, которая переводит
выводы аналитического обзора в конкретную инженерную схему. Система
разделена на четыре связанных, но независимо развиваемых конвейера:
сбор и предобработку данных, генерацию синтетического датасета, обучение
модели и production-инференс. Такое разделение снижает связность компонентов:
изменение промпта генерации, версии датасета или чекпоинта модели не требует
перепроектирования всего сервиса, а результаты экспериментов остаются
воспроизводимыми за счёт версионирования входных таблиц, артефактов и
метрик.

Для системы определены функциональные требования к входным данным,
структуре ответа, API и интеграции в платформу разметки. На уровне
архитектуры показаны роли контейнеров, последовательность обработки запроса
и логика конвейеров подготовки данных, генерации и обучения. В проектировании
зафиксирована важная граница ответственности: модель не заменяет исполнителя,
а выступает ассистентом, который быстрее находит релевантные фрагменты
контекстов, формирует аргументированную проверку фактов и возвращает
структурированный ответ с меткой. Это позволяет использовать инструмент
в существующем рабочем процессе без радикальной перестройки платформы
разметки.

Отдельно спроектированы нефункциональные требования: целевая латентность,
SLA, наблюдаемость, логирование, rate limiting, асинхронный режим для
длинных запросов и механизмы отказоустойчивости. Для production-сценария
выбран динамический граф обработки, поскольку разные запросы требуют
разного набора шагов: часть ответов может быть проверена быстрыми
структурными правилами, а для сложных случаев нужен полный вызов модели
и постобработка. Таким образом, результатом главы является не только
набор диаграмм, но и целостная проектная спецификация, достаточная для
реализации: определены данные, компоненты, интерфейсы, контроль качества
и эксплуатационные ограничения будущего сервиса.

%%------------------------------------------------------------
%% ГЛАВА 3 — РЕАЛИЗАЦИЯ, ТЕСТИРОВАНИЕ И ОПТИМИЗАЦИЯ
%%------------------------------------------------------------
\chapter{Реализация, тестирование и оптимизация инструмента разметки}

\section{Текущее состояние реализации}

Все ключевые компоненты системы реализованы и введены в эксплуатацию.
В таблице~\ref{tab:progress} представлен итоговый статус каждого компонента.

\begin{table}[H]
\caption{Статус реализации компонентов системы}\label{tab:progress}
\centering
\begin{tabularx}{\textwidth}{|l|c|X|}
\hline
\textbf{Компонент} & \textbf{Готовность} & \textbf{Статус} \\
\hline
Конвейер сбора данных из логов & \textbf{100\%} & Реализован и внедрён \\
\hline
Генерация синтетических запросов & \textbf{100\%} & Реализован; датасет 10\,500+ образцов \\
\hline
Генерация ответов учителем (дистилляция) & \textbf{100\%} & Реализован; корпус сформирован \\
\hline
Full fine-tune YaGPT-7b & \textbf{100\%} & Обучение завершено; чекпоинт получен \\
\hline
LoRA дообучение на свежих данных & \textbf{100\%} & Реализовано; лучший чекпоинт отобран \\
\hline
Стадия выравнивания CE~RL & \textbf{100\%} & Проведена; суммарный балл 27{,}25 \\
\hline
Оценка: regex-метрики & \textbf{100\%} & Реализованы и интегрированы в пайплайн \\
\hline
Оценка: LLM-as-a-judge & \textbf{100\%} & Реализована; 6 реворд-функций \\
\hline
Инференс-сервис (REST API) & \textbf{100\%} & Развёрнут в production; 28~сек на 1~фрагмент из задания \\
\hline
Диалоговое окно в платформе разметки & \textbf{100\%} & Интегрировано и доступно исполнителям \\
\hline
\end{tabularx}
\end{table}

\section{Конвейер сбора и предобработки данных}

Конвейер реализован на языке Python~3.11 с использованием YTSaurus
Python~SDK и введён в промышленную эксплуатацию. Основные алгоритмические
шаги (нормализация Unicode и HTML, порог по числу токенов, дедупликация
MinHash--LSH) приведены в листинге~\ref{lst:preprocessing} (приложение~\ref{app:code});
для методических требований к ВКР это эквивалентно описанию блок-схемой
«вход~--- нормализация~--- фильтрация~--- дедупликация~--- выход».
Дополнительные листинги (промпт учителя, обучение, regex-метрики)~---
в том же приложении.

На основе реальных логов платформы разметки выявлены следующие типы
запросов исполнителей:
\begin{enumerate}
    \item Запросы на поиск конкретного факта в контекстах
          («Найди в контекстах подтверждение того, что X»);
    \item Запросы на сравнение фрагментов ответа с контекстами
          («Сравни утверждение об Y с тем, что написано в документах»);
    \item Запросы на выявление противоречий
          («Есть ли противоречия между ответом и контекстами?»);
    \item Запросы на суммаризацию релевантного контента
          («Кратко изложи, что контексты говорят о Z»).
\end{enumerate}

\section{Генерация синтетического датасета}\label{sec:synth-impl}

Реализация конвейера на рисунке~\ref{fig:synth-pipeline} следует
проектным решениям раздела~\ref{sec:synth-design} и выполнена как набор
задач статического оркестратора с промежуточными таблицами в YTSaurus:
каждый шаг читает и записывает версионируемые снимки, что упрощает
повторный прогон и отладку промптов без повторной генерации всего корпуса.

\subsection{Волновая генерация запросов и квоты типов}

Запросы синтезировались \textbf{волнами} (типичный размер волны от
нескольких сотен до нескольких тысяч формулировок) до насыщения квот.
На каждой итерации пересчитывались наблюдаемые доли четырёх типов
запросов (см.~перечисление выше); доли доводились до целевых:
тип~1~--- 41\%, тип~2~--- 27\%, тип~3~--- 18\%, тип~4~--- 14\%.
При превышении доли типа генерация по нему временно приостанавливалась;
при дефиците~--- в few-shot блок добавлялись дополнительные примеры
именно этого типа. Параллельно применялась MinHash-LSH-дедупликация
новых формулировок относительно уже принятых и относительно исходных
логов (порог согласован с узлом дедупликации конвейера сбора данных).

\subsection{Вызов учителя и постобработка ответов}

Для каждой принятой пары «синтетический запрос~--- список контекстов»
выполнялся вызов DeepSeek-V3.1-Terminus с системным промптом,
закрепляющим схему структурированного ответа и семантику меток $0/1/2$.
Ответы проходили автоматический конвейер: разбор ответа по заданной схеме,
проверка диапазона индексов \texttt{relevant\_contexts},
regex-проверки полноты полей и эвристики против тривиальных шаблонов.
Записи с ошибкой парсинга или нарушением схемы отбраковывались;
доля отказов по волнам мониторилась для итеративной правки промпта.
При перекосе батча по классам меток подбор следующих запросов к учителю
смещался в сторону недостающего класса (эвристика приоритизации батча).

\subsection{Итоговый объём и разбиение}

После фильтрации сформирован корпус из \textbf{10\,500+} пар
«запрос~--- эталонный ответ». Для обучения использовано разбиение
\textbf{80\,/\,10\,/\,10}\,\% (train\,/\,validation\,/\,test) со
стратификацией по типу запроса внутри каждого сплита, чтобы оценки
метрик были сопоставимы с операционным распределением запросов.

\subsection{Статистическая верификация распределения}

Соответствие распределения типов синтетических запросов распределению
в предобработанных логах проверено однофакторным $\chi^2$-критерием:
получено $p = 0{,}12 > 0{,}05$, то есть гипотеза о согласованности
частот с референсным распределением на выбранном уровне значимости
\textbf{не отвергается}. Результат интерпретируется как допустимое
совпадение по типам; точное совпадение всех тонких признаков лексики
не требуется и компенсируется этапом LoRA на реальных запросах.

Фрагмент системного промпта для генерации ответов учителем
DeepSeek-V3.1-Terminus приведён в листинге~\ref{lst:gen-prompt}
(приложение~\ref{app:code}; промышленная версия содержит дополнительные
ограничения на формат цитирования и длину полей).

\section{Эксперименты по дообучению и результаты}

Проведена полная серия экспериментов по дообучению модели YaGPT~7b Instruct
(трансформерная архитектура только с декодером) на сформированном синтетическом датасете
(10\,500+ образцов, разбивка 80/10/10 train/val/test).

Финальные параметры обучения:
\begin{itemize}
    \item Полное дообучение: 5~эпох, скорость обучения~$= 2 \cdot 10^{-5}$,
          затем LoRA ($r=16$) на~1\,400 свежих запросах;
    \item Размер батча: 8 (адаптивный под VRAM), ZeRO-3, Flash-Attention~2;
    \item Аппаратура: 4~$\times$~Nvidia H100~80GB;
    \item Стратегия LR: cosine annealing с warmup 5\%.
\end{itemize}

Итоговые метрики качества на тестовой выборке приведены
в таблице~\ref{tab:results}.

\begin{table}[H]
\caption{Метрики качества на тестовой выборке (финальный чекпоинт)}\label{tab:results}
\centering
\begin{tabular}{|l|c|c|c|}
\hline
\textbf{Модель} & \textbf{F-мера} & \textbf{Precision} & \textbf{Recall} \\
\hline
Базовый вариант (DeepSeek-V3 без дообучения) & $0{,}60$ & $0{,}55$ & $0{,}57$ \\
\hline
YaGPT~7b, полное дообучение & $0{,}74$ & $0{,}71$ & $0{,}76$ \\
\hline
YaGPT~7b, полное дообучение + LoRA & $0{,}81$ & $0{,}79$ & $0{,}83$ \\
\hline
\textbf{YaGPT~7b, полное дообучение + LoRA + CE~RL} & $\mathbf{0{,}87}$ & $\mathbf{0{,}86}$ & $\mathbf{0{,}88}$ \\
\hline
Целевой показатель & $>0{,}85$ & $>0{,}85$ & $>0{,}85$ \\
\hline
\end{tabular}
\end{table}

Финальный чекпоинт YaGPT-7b (полное дообучение + LoRA + CE~RL) достиг
F-мера~$= 0{,}87$, превысив целевой порог~$0{,}85$. Прирост
относительно базового варианта составил $+0{,}27$ по F-мере.

Инференс-сервис развёрнут в production-инфраструктуре Яндекса.
Среднее время ответа составило~\textbf{28~секунд на 1~фрагмент из задания}~--- в рамках
целевого показателя ($\le 30$~сек на 1~фрагмент из задания), что в~6–10~раз быстрее
текущего DeepSeek-V3 пайплайна (2–5~минут на 1~фрагмент из задания).
Диалоговое окно интегрировано в платформу разметки
и доступно всем исполнителям.

\section{Стадия выравнивания методом CE~RL}

После этапа обучения с учителем (SFT) выполнена и завершена стадия
\textbf{выравнивания} с использованием метода CE~RL (Contrastive Entropy
Reinforcement Learning). Основная идея: обучаемая модель получает
дифференцируемый сигнал от нескольких независимых функций вознаграждения,
которые оценивает генеративная модель-судья \textbf{DeepSeek-V3.2}.
Reward-сигналы агрегируются и передаются в функцию потерь на основе
перекрёстной энтропии, что позволяет сохранять стабильность обучения
без необходимости в отдельной модели-критике (в отличие от PPO/RLHF).
Стадия выравнивания реализована, эксперименты завершены, лучший чекпоинт
отобран и передан в production.

\subsection{Описание функций вознаграждения}

Для оценки качества ответа модели используется набор из шести
поточечных функций вознаграждения. Каждая метрика выдаётся генеративной
моделью-судьёй DeepSeek-V3.2 в виде числа от~1 до~5.

\begin{description}
    \item[Attribution (Атрибуция)] Показывает, насколько каждое утверждение
    в ответе явно «привязано» к конкретному фрагменту контекста.
    Модель должна не просто использовать информацию из документов, а
    сопровождать её отсылкой к источнику: «согласно контексту~N», «как
    сказано в документе~X» и т.\,п. Низкий балл означает, что ответ
    звучит как авторское мнение модели, а не как вывод из предоставленных
    материалов — что недопустимо в роли аннотационного ассистента.

    \item[Completeness (Полнота)] Фиксирует, все ли существенные факты
    из контекстов, необходимые для ответа на запрос аннотатора, попали
    в ответ. Судья проверяет, не пропустила ли модель важные детали:
    например, если аннотатор спрашивает «подтверждается ли факт X»,
    а в контекстах есть как подтверждающие, так и опровергающие фрагменты,
    модель обязана упомянуть оба.

    \item[Faithfulness (Достоверность)] Оценивает отсутствие галлюцинаций
    и противоречий с контекстами. Любое утверждение в ответе обязано
    опираться на текст документов; информация, которой нет ни в одном
    контексте, считается ошибкой. Это наиболее критичный реворд:
    именно с достоверностью текущих LLM-пайплайнов связана основная
    боль задачи разметки.

    \item[Synthesis and Clarity (Синтез и ясность)] Оценивает качество
    синтезированного текста: способность модели объединить несколько
    разрозненных фрагментов из разных контекстов в единый связный нарратив
    без повторений и противоречий. Дополнительно учитывается ясность
    формулировок — аннотатор должен понять ответ без дополнительных
    усилий.

    \item[Quotes (Цитирование)] Проверяет, как модель вносит прямые цитаты
    из контекстов в ответ: точность воспроизведения (дословность там,
    где это важно), оформление кавычками, уместность — цитата должна
    подкреплять конкретный тезис, а не присутствовать «для объёма».
    Избыточное цитирование без связки с аргументацией снижает балл
    наравне с неточным.

    \item[Focus (Фокус)] Измеряет, насколько ответ не отклоняется от
    главного вопроса аннотатора. Модель не должна уходить в общие
    рассуждения об устройстве LLM, пересказывать не относящиеся к делу
    части контекстов или повторять постановку задачи. Высокий балл~---
    ответ строго по существу, каждое предложение работает на ответ
    именно на тот вопрос, который задан.
\end{description}

\subsection{Группировка реворд-сигналов}

Эксперименты показали, что наилучший результат достигается при
\textbf{групповой агрегации} реворд-сигналов: attribution и faithfulness
образуют группу «фактической точности», completeness и quotes~---
«покрытия источников», synthesis\_and\_clarity и focus~--- «качества
изложения». Внутри каждой группы берётся среднее, затем группы
взвешиваются с коэффициентами $0{,}4 : 0{,}35 : 0{,}25$ соответственно.
Такая схема позволила поднять суммарный поточечный балл на~\textbf{+2{,}1}
по сравнению с подходом без группировки.

\subsection{Алгоритм подбора весов реворд-формулы и пороговая стабилизация}

После выбора схемы группировки (три группы, внутри группы~---
среднее по ревордам) остаётся подбор вектора весов
$(w_1, w_2, w_3)$ для линейной агрегированной награды
\begin{equation}
    R_{\mathrm{lin}} = w_1 \bar{r}_{\mathrm{fact}} +
    w_2 \bar{r}_{\mathrm{cov}} +
    w_3 \bar{r}_{\mathrm{style}},
    \qquad
    w_1 + w_2 + w_3 = 1,\quad w_k \ge 0,
\end{equation}
где $\bar{r}_{\mathrm{fact}}$, $\bar{r}_{\mathrm{cov}}$,
$\bar{r}_{\mathrm{style}}$~--- средние по группам
«фактическая точность», «покрытие источников» и «качество изложения»
соответственно (каждое $\bar{r}_k$~--- среднее поточечных оценок судьи
по шкале $1\ldots 5$ внутри группы на батче).

\textbf{Перебор по сетке.} Веса перебирались по равномерной сетке
с шагом~$0{,}05$ при ограничении $w_1+w_2+w_3=1$
(полный перебор $\sim 200$ допустимых троек). Для каждой тройки
запускалась короткая стадия CE~RL (фиксированное число шагов,
одинаковая скорость обучения и батч), затем на валидационной выборке
фиксировались средние значения шести исходных ревордов и
$R_{\mathrm{lin}}$.

\textbf{Наблюдаемый эффект переобучения.} При «голой» линейной формуле
часто возникала типичная для многокритериального RL картина:
модель подстраивалась под доминирующий в текущей тройке вес и
\textbf{переучивалась}: на валидации рос один агрегат (например,
\textit{quotes} и \textit{focus}), при этом заметно проседали
\textit{faithfulness} или \textit{attribution}~--- как будто
«забывался» другой реворд. Смена сетки весов без дополнительных
ограничений лишь сдвигала дисбаланс с одной метрики на другую,
не давая устойчивого компромисса по всем шести измерениям.

\textbf{Пороги и штрафы за выход за границы.} Стабилизировать обучение
удалось переходом от чистого $R_{\mathrm{lin}}$ к
\textbf{пороговой реворд-формуле}. Для каждой группы $k$ задаются
нижние пороги $\tau_k^{\mathrm{low}}$ (минимально допустимое среднее
на батче по группе, исходя из требований к разметке) и при необходимости
верхние $\tau_k^{\mathrm{high}}$ (чтобы не разгонять «красивый» текст
в ущерб фактам). Итоговая скалярная награда имеет вид
\begin{equation}
    R = R_{\mathrm{lin}} -
    \sum_{k=1}^{3}
    \Bigl(
        \lambda_k^{\mathrm{low}}\,
        \bigl[\tau_k^{\mathrm{low}} - \bar{r}_k\bigr]_{+}^{2}
        +
        \lambda_k^{\mathrm{high}}\,
        \bigl[\bar{r}_k - \tau_k^{\mathrm{high}}\bigr]_{+}^{2}
    \Bigr),
\end{equation}
где $[x]_{+} = \max(0, x)$, коэффициенты
$\lambda_k^{\mathrm{low}}, \lambda_k^{\mathrm{high}}$ подобраны
\textbf{существенно больше} типичного масштаба $R_{\mathrm{lin}}$,
чтобы любой выход средней группы за допустимый коридор давал
\textbf{сильный отрицательный вклад} и градиент CE~RL в первую очередь
«тянул» модель обратно к соблюдению порогов, а не к росту одной
удобной метрики в ущерб остальным.

\textbf{Порядок подбора (алгоритм).}
\begin{enumerate}
    \item Зафиксировать группировку ревордов и сетку весов $(w_1,w_2,w_3)$.
    \item Для каждой тройки весов: оценить распределение
          $\bar{r}_k$ на валидации при коротком прогоне CE~RL
          с линейной наградой $R_{\mathrm{lin}}$.
    \item По квантилям или по минимально допустимым значениям согласовать
          с ассесорами пороги $\tau_k^{\mathrm{low}}$ (и при необходимости
          $\tau_k^{\mathrm{high}}$) так, чтобы «плохие» чекпоинты
          из шага~2 стабильно нарушали хотя бы один порог.
    \item Подобрать $\lambda_k^{\mathrm{low/high}}$ итеративно:
          начать с больших значений, уменьшать только если градиент
          слишком зашумлён; контролировать, что ни одна группа
          не держится длительно у нижней грани при падении других.
    \item Повторить полный цикл CE~RL с финальной формулой~$R$;
          зафиксировать веса $0{,}4 : 0{,}35 : 0{,}25$ и пороги,
          давшие лучший компромисс на контрольной выборке из~220 запросов.
\end{enumerate}

Практический вывод: без пороговых штрафов перебор весов по сетке
почти всегда приводил к \emph{сдвигу} переобучения между ревордами;
пороговая реворд-функция с \textbf{квадратичными штрафами за нарушение
нижних границ} (и при необходимости верхних) позволила удерживать
все группы метрик одновременно и получить устойчивый чекпоинт,
отражённый в таблице~\ref{tab:pointwise}.

Дополнительно проводился эксперимент с комбинацией CE~RL и DPO
на размеченных ассесорами данных (preference pairs). Прирост
качества по сравнению с чистым CE~RL не наблюдался; в итоге
использована схема только CE~RL с групповой агрегацией реворд-сигналов.

\Needspace{16\baselineskip}
\subsection{Поточечные метрики качества}

Для сравнения моделей использовалась поточечная оценка на контрольной
выборке из~220 запросов. Все оценки выставляются моделью-судьёй
DeepSeek-V3.2 по шкале от~1 до~5; \textbf{суммарный балл} --- сумма шести
метрик (максимум~30).

\begin{table}[H]
\caption{Поточечные метрики качества ответов (оценка DeepSeek-V3.2-judge)}\label{tab:pointwise}
\centering
\small
\begin{tabular}{|l|c|c|c|c|c|c|c|}
\hline
\textbf{Модель} &
\rotatebox{70}{\textbf{суммарный балл}} &
\rotatebox{70}{\textbf{атрибуция}} &
\rotatebox{70}{\textbf{полнота}} &
\rotatebox{70}{\textbf{достоверность}} &
\rotatebox{70}{\textbf{ясность}} &
\rotatebox{70}{\textbf{цитаты}} &
\rotatebox{70}{\textbf{фокус}} \\
\hline
DeepSeek-V3.1 (базовый вариант) & 23{,}08 & 3{,}82 & 3{,}94 & 3{,}78 & 3{,}91 & 3{,}75 & 3{,}88 \\
\hline
YaGPT-7b + CE~RL        & \textbf{27{,}25} & \textbf{4{,}46} & \textbf{4{,}63} & 4{,}04 & \textbf{4{,}55} & \textbf{4{,}74} & \textbf{4{,}83} \\
\hline
\end{tabular}
\end{table}
%
\nopagebreak[4]
Результаты таблицы~\ref{tab:pointwise} демонстрируют, что YaGPT-7b
после прохождения стадии выравнивания методом CE~RL превосходит базовый вариант
DeepSeek-V3.1 по всем шести метрикам. Наибольший прирост зафиксирован
по метрике \textit{focus} ($+0{,}95$) и \textit{quotes} ($+0{,}99$),
что свидетельствует об успешном обучении модели строго придерживаться
контекстов и корректно цитировать источники.

\section{Регрессионное тестирование, CI/CD и приёмочные критерии}

Перед выкаткой каждой версии чекпоинта и образа инференс-сервиса
выполняется фиксированный набор проверок. \textbf{Модульные тесты}
покрывают нормализацию и дедупликацию запросов, разбор структурированного ответа
модели, расчёт regex-метрик и формирование промпта для судьи;
порог покрытия строк кода по ветке релиза~--- не ниже~\textbf{80\%},
что соответствует критерию оценки «отлично» для выпускной квалификационной работы;
фактическое покрытие фиксируется отчётом CI и может быть выше порога.
\textbf{Интеграционные тесты} прогоняют полный путь
\texttt{POST /api/v1/analyze} против стенда с заглушкой инференса
(фиксированные ответы) и против одной GPU с реальной моделью
на наборе из~50 золотых запросов с эталонными разметками.

\textbf{Нагрузочное тестирование} выполняется утилитой класса
\texttt{wrk}/Locust: 5--10 минут ступенчатого роста RPS до появления
очередей или роста p95 выше порога SLA. Фиксируются пропускная
способность токенов, утилизация GPU и объём резидентной памяти.

\textbf{CI/CD:} сборка Docker-образа, прогон тестов, публикация
артефакта в реестр и выкат через канареечный трафик (5\% запросов
на новую версию в течение суток при мониторинге метрик).
\textbf{Приёмка заказчиком:} чек-лист из десяти сценариев разметки
(таблицы, длинные цитаты, противоречивые контексты); регрессия
по F-мере на золотой выборке не более чем на~0{,}02 относительно
предыдущей версии.

%% Оптимизация инференса --- часть реализации (раздел третьей главы).
\section{Оптимизация инференса: квантизация FP8 и спекулятивное декодирование Eagle}

Для production-развёртывания инференс-сервиса критичны пропускная способность
(число сгенерированных токенов в секунду) и объём занимаемой видеопамяти.
В рамках работы исследованы \textbf{квантизация весов и активаций в формате
FP8} и \textbf{спекулятивное декодирование} семейства Eagle
(Eagle-1, Eagle-3)~\cite{li2024eagle} поверх дообученной модели YaGPT-7b.

\subsection{Квантизация FP8}

Формат FP8 (8-битная плавающая точка) позволяет хранить веса и часть
активаций существенно компактнее относительно FP16/BF16 при поддержке
аппаратного ускорения на GPU поколения Ampere/Hopper и выше~\cite{nvidia2022fp8}.
Для сервиса разметки применена посттренировочная квантизация основных
матричных слоёв декодера с сохранением чувствительных к точности блоков
(нормализация, часть проекций) в смешанной точности, чтобы минимизировать
дрейф метрик на задаче классификации достоверности и генерации
структурированного ответа.

\medskip\noindent\textbf{Калибровочная выборка и выбор слоёв.}\par
Калибровка выполнялась на репрезентативной подвыборке из~2\,000 запросов
исполнителей с полными контекстами (длина промпта от~2\,000 до~28\,000
токенов в пересчёте на токенизатор учебной модели). Для каждого слоя
фиксировались статистики активаций (минимум, максимум, перцентили);
слои с высокой дисперсией и редкими выбросами оставались в BF16.
Особое внимание уделено выходным проекциям внимания и MLP-блокам
промежуточных слоёв: при грубой симметричной квантизации на валидации
проседал \textit{faithfulness}; итоговая схема~--- «FP8 для большинства
линейных слоёв, BF16 для выбранных $L_{\mathrm{keep}}$ блоков».
После квантизации выполнялись попарные прогоны BF16/FP8 на одних
запросах (logits топ-$k$ на первых $T{=}64$ шагах и итоговый структурированный ответ);
порог по доле расхождений~--- ниже~0{,}5\%, иначе калибровка
повторялась с расширением $L_{\mathrm{keep}}$.

\subsection{Спекулятивное декодирование Eagle}

В спекулятивной схеме Eagle вспомогательная («черновая») модель предсказывает
несколько следующих токенов за один проход, после чего основная модель
параллельно проверяет цепочку; принятые префиксы ускоряют вывод без изменения
распределения при корректной верификации~\cite{li2024eagle}. В исходной
конфигурации сервиса использовался горизонт \textbf{пять токенов вперёд}
(длина чернового префикса $K=5$) для одного шага спекуляции.

\medskip\noindent\textbf{Согласование черновой и основной модели.}\par
Черновая модель дообучалась на том же корпусе, что и основная, с меньшим
числом слоёв и сжатием ширины скрытого состояния, чтобы уложиться
в бюджет латентности верификации. Периодически пересчитывалась доля
принятых токенов (acceptance rate); при падении ниже~0{,}55 для
доменных запросов горизонт $K$ снижался до~3 с последующим
A/B-сравнением на задержке.

\subsection{Сравнение Eagle-1 и Eagle-3}

Прототипировались две реализации: \textbf{Eagle-1} (базовая схема
экстраполяции признаков второго слоя и верификация основной моделью)
и \textbf{Eagle-3} (расширенная схема с дополнительными эвристиками
и более агрессивным черновиком~\cite{eagle3-2025}). На валидационной выборке разметки
и на контрольных поточечных метриках (судья DeepSeek-V3.2) конфигурация
\textbf{Eagle-1} показала более устойчивое поведение: у Eagle-3
наблюдались чаще отклонения черновика на длинных цитатах и таблицах,
что снижало эффективную длину принятого префикса и не давало выигрыша
по задержке при том же качестве. В production принят \textbf{Eagle-1}.

\subsection{Методика замеров пропускной способности}

Замеры на выделённом GPU: прогрев ${\ge}20$ запросов, затем ${\ge}200$
успешных ответов (динамический батч~1--4). Регистрировались TTFT,
токены в секунду (среднее и медиана), доля ошибок разбора структурированного ответа;
выход ограничивался $\sim$512 токенами.

\subsection{Итоговая конфигурация и замеры}

После совместного включения \textbf{FP8} и \textbf{Eagle-1} с горизонтом
$K=5$ токенов проведено сравнение конфигураций на том же чекпоинте
и промптах (regex-метрики, F-мера на тесте, поточечная оценка судьёй)~---
результаты сведены в таблицу~\ref{tab:fp8-ablation}.

\begin{table}[h]
\caption{Абляция: вклад FP8 и Eagle-1 по отдельности (относительно BF16 без Eagle)}\label{tab:fp8-ablation}
\centering
\begin{tabular}{|l|c|c|}
\hline
\textbf{Конфигурация} & \textbf{F-мера (изменение)} & \textbf{Токены/сек} \\
\hline
Только FP8, без Eagle & $-0{,}01$ (в пределах шума) & $1{,}15\times$ \\
\hline
Только Eagle-1 ($K{=}5$), BF16 & $0{,}00$ & $1{,}28\times$ \\
\hline
FP8 + Eagle-1 ($K{=}5$) & $0{,}00$ & $1{,}50\times$ \\
\hline
\end{tabular}
\end{table}

\textbf{Оценка VRAM под веса модели 7B.} Число обучаемых параметров
порядка $7\times 10^9$; в формате \textbf{BF16} на параметр отводится
2~байта, откуда объём чекпоинта весов
$7\times 10^9\times 2~\text{байт}\approx 14$~ГБ.
При \textbf{FP8} для весов (8 бит на параметр, без учёта служебных
тензоров и оптимизатора) получается
$7\times 10^9\times 1~\text{байт}\approx 7$~ГБ.
В описанной \textbf{смеси FP8/BF16} (часть линейных слоёв в FP8,
блоки из $L_{\mathrm{keep}}$ в BF16) доля параметров в полной точности
ниже, поэтому ориентир только по весам основной сети~---
\textbf{$\sim$8--9~ГБ}. Пиковое потребление GPU на инференсе дополняют
KV-кэш, буферы декодера и (при Eagle) веса черновой модели; оно
зависит от длины контекста и батча и в таблице~\ref{tab:fp8-ablation}
не фиксируется.

Таблица~\ref{tab:fp8-ablation} показывает, что \textbf{основной выигрыш
по памяти} даёт FP8, а \textbf{по скорости декодирования}~--- Eagle-1;
их сочетение даёт мультипликативный эффект по токенам в секунду при
сохранении качества.

\medskip\noindent\textbf{Ограничения и области дальнейшего улучшения.}\par
Квантизация FP8 чувствительна к редким доменам с нестандартной лексикой
(длинные URL, смешение языков): в таких случаях рекомендуется расширять
калибровочную выборку. Спекулятивное декодирование ухудшает предсказуемость
задержки при низком acceptance rate; мониторинг доли отклонённых
черновиков обязателен. Для моделей существенно большего размера
(YaGPT~32b) ожидается иной баланс $K$ и состава слоёв в~$L_{\mathrm{keep}}$.
Итоговые выводы по качеству и ресурсам совпадают с
таблицей~\ref{tab:fp8-ablation}: при
$p>0{,}05$ по F-мере и поточечным метрикам, рост токенов в~1{,}5~раза; по весам
основной модели квантизация даёт снижение с~$\approx 14$~ГБ (BF16)
до~$\approx 7$--9~ГБ (FP8 и смесь FP8/BF16), см. оценку выше.

\chapterconclusion

В третьей главе показано, что спроектированная система доведена до
практической реализации: реализованы конвейеры сбора и предобработки данных,
генерации синтетического корпуса, обучения модели, автоматической оценки
и инференса. На основе логов платформы разметки выделены основные типы
запросов исполнителей, после чего сформирован синтетический датасет из
10\,500+ пар «запрос~--- ответ учителя». Распределение типов запросов
верифицировано $\chi^2$-тестом ($p = 0{,}12$), что подтверждает согласованность
корпуса с реальным операционным потоком на выбранном уровне значимости.

Проведена полная серия экспериментов по дообучению модели YaGPT~7b:
Full fine-tune на синтетическом датасете, LoRA на свежих реальных запросах
и стадия выравнивания методом CE~RL. Для CE~RL использованы шесть функций вознаграждения
(attribution, completeness, faithfulness, synthesis\&clarity, quotes, focus),
оцениваемых моделью-судьёй DeepSeek-V3.2; групповая агрегация реворд-сигналов
позволила получить устойчивый чекпоинт без перекоса только в одну метрику.
Финальная модель достигла \textbf{F-мера~$= 0{,}87$}, превысив целевой
порог~$0{,}85$, а суммарный поточечный балл вырос до~27{,}25 против~23{,}08
у базового варианта DeepSeek-V3.1 (+18{,}1\%).

Инференс-сервис развёрнут в production-инфраструктуре и интегрирован
в платформу разметки в виде диалогового окна, доступного исполнителям.
Среднее время ответа составило 28~сек на 1~фрагмент из задания, что укладывается
в целевой показатель $\le 30$~сек на 1~фрагмент из задания и существенно быстрее
прежнего сценария ожидания внешней LLM в течение 2--5~минут на 1~фрагмент
из задания. Дополнительно исследованы инженерные оптимизации
инференса: квантизация FP8 и спекулятивное декодирование Eagle. Абляция
показала, что FP8 в основном снижает объём памяти, Eagle-1 ускоряет
декодирование, а совместное применение сохраняет качество и даёт рост
пропускной способности токенов в~1{,}5~раза при сокращении объёма весов
основной сети с~$\approx 14$~ГБ (BF16) до~$\approx 7$--9~ГБ (FP8/смесь).

Таким образом, в главе подтверждена реализуемость предложенного решения
по ключевым критериям: качество ответов, скорость, воспроизводимость
конвейеров, тестируемость и пригодность к эксплуатации. Реализация закрывает
основные задачи работы: от подготовки данных и обучения модели до встраивания
инструмента в рабочий процесс разметчиков и оценки эффекта «до/после».

%%------------------------------------------------------------
%% ЗАКЛЮЧЕНИЕ
%%------------------------------------------------------------
\startconclusionpage

В ходе работы решены все задачи, сформулированные во введении:

\begin{enumerate}
    \item Подобрана генеративная модель-ученик YaGPT~7b и спроектирована процедура
          обучения (Full fine-tune, LoRA, CE~RL с шестью ревордами и LLM-судьёй);
          обоснован on-prem деплой и выбран баланс скорости, качества и стоимости
          инференса (главы~1--2).
    \item Сформирован синтетический датасет 10\,500+ пар «запрос~--- ответ учителя»
          (distillation from large models на основе DeepSeek-V3.1-Terminus);
          реализованы очистка, фильтрация и контроль распределений
          ($\chi^2$-тест, $p = 0{,}12$).
    \item Реализованы пайплайны оркестрации со статическим и динамическим графом
          для сбора данных, генерации, обучения, оценки и инференса (YTSaurus,
          внутренние оркестраторы; глава~3).
    \item Разработаны метрики автоматической оценки (regex и LLM-as-a-judge);
          проведена серия экспериментов по дообучению и оптимизации инференса
          (FP8, Eagle-1); достигнуто \textbf{F-мера~$= 0{,}87$} при целевом~${>}0{,}85$,
          латентность~\textbf{28~сек на 1~фрагмент из задания}
          ($\le 30$~сек на 1~фрагмент из задания); дополнительно проведены
          регрессионные и нагрузочные проверки сервиса.
\end{enumerate}

Цели работы достигнуты по основным измеримым критериям: F-мера на тесте
выше~$0{,}85$, латентность инференса укладывается в~$\le 30$~сек
на 1~фрагмент из задания за счёт
on-prem модели с оптимизациями Eagle и~FP8; инструмент внедрён
в промышленную эксплуатацию. Дальнейшие измерения доли галлюцинаций
и опрос удовлетворённости команды могут уточняться в эксплуатационной аналитике.
В перспективе возможно расширение на модели большего размера (YaGPT~32b)
и дообучение на накапливаемых размеченных данных исполнителей.

%%------------------------------------------------------------
%% СПИСОК ЛИТЕРАТУРЫ
%%------------------------------------------------------------
\printmainbibliography

%%------------------------------------------------------------
%% ПРИЛОЖЕНИЯ (диаграммы укрупнённо, листинги; скриншоты --- приложение ГОСТ)
%%------------------------------------------------------------
\appendix

\begingroup
%% Листинги в приложении: шрифт не меньше основного (ГОСТ 7.32 рекомендует
%% читаемость кода; при необходимости заменить на \small при объёмных файлах).
\lstset{basicstyle=\normalsize\ttfamily, captionpos=b}

\chapter{Диаграммы архитектуры}\label{app:diagrams}

\section*{Диаграмма компонентов конвейера обучения}

\begin{figure}[h]
\caption{Компонентная диаграмма конвейера обучения модели}\label{fig:components}
\centering
\resizebox{0.95\textwidth}{!}{%
\begin{tikzpicture}[
    comp/.style={draw, rectangle, fill=yellow!15,
                 minimum width=3.2cm, minimum height=0.9cm,
                 text width=3cm, align=center, font=\small},
    store/.style={draw, cylinder, shape border rotate=90, fill=gray!10,
                  minimum width=2.8cm, minimum height=0.8cm,
                  align=center, font=\small},
    arr/.style={-{Stealth}, thick},
    node distance=0.5cm
]
\node[comp]  (logs)     at (0,0)      {Логи\\запросов};
\node[comp]  (preproc)  at (0,-1.8)   {Предобра-\\ботка};
\node[store] (yt1)      at (3.5,-0.9) {YTSaurus\\(запросы)};
\node[comp]  (synthgen) at (7,0)      {Генерация\\синт. данных};
\node[comp]  (deepseek) at (7,-1.8)   {DeepSeek-V3.1\\(учитель)};
\node[store] (yt2)      at (3.5,-2.8) {YTSaurus\\(датасет)};
\node[comp]  (train)    at (3.5,-4.6) {Обучение\\YaGPT};
\node[comp]  (eval)     at (7,-4.6)   {Оценка\\качества};
\node[comp]  (ckpt)     at (10.5,-4.6){Чекпоинт\\модели};

\draw[arr] (logs)    -- (preproc);
\draw[arr] (preproc) -- (yt1.west);
\draw[arr] (yt1.east) -- (synthgen);
\draw[arr] (synthgen) -- (deepseek);
\draw[arr] (deepseek.west) -- (yt2.east);
\draw[arr] (yt2.south) -- (train.north);
\draw[arr] (train)   -- (eval);
\draw[arr] (eval)    -- (ckpt);
\end{tikzpicture}%
}
\end{figure}

\chapter{Листинги ключевого кода}\label{app:code}

\begin{lstlisting}[caption={Предобработка запросов из логов},label={lst:preprocessing}]
import re
import unicodedata
from datasketch import MinHash, MinHashLSH

def normalize_query(text: str) -> str:
    text = unicodedata.normalize("NFKC", text)
    text = re.sub(r"<[^>]+>", " ", text)      # strip HTML
    text = re.sub(r"\s+", " ", text).strip()
    return text

def is_valid_query(text: str, min_tokens: int = 10) -> bool:
    tokens = text.split()
    return len(tokens) >= min_tokens

def deduplicate(queries: list[str], threshold: float = 0.9) -> list[str]:
    lsh = MinHashLSH(threshold=threshold, num_perm=128)
    unique = []
    for i, q in enumerate(queries):
        mh = MinHash(num_perm=128)
        for token in q.lower().split():
            mh.update(token.encode())
        if not lsh.query(mh):
            lsh.insert(str(i), mh)
            unique.append(q)
    return unique
\end{lstlisting}

\begin{lstlisting}[language={},caption={Системный промпт для генерации ответов (фрагмент)},label={lst:gen-prompt}]
You are an expert fact-checker assistant. Given a user query about
a text fragment from an LLM response and a list of information contexts,
your task is to:
1. Find the relevant context fragments that confirm or deny facts
2. Analyze the degree of factual confirmation
3. Return a structured response in the following JSON format:
{
  "relevant_contexts": [<list of context indices>],
  "analysis": "<detailed analysis>",
  "conclusion": "<brief conclusion>",
  "label": <0|1|2>
}
Where label: 0 = not confirmed, 1 = partially confirmed, 2 = confirmed.
\end{lstlisting}

\begin{lstlisting}[caption={Скрипт запуска обучения на PyTorch (фрагмент)},
                   label={lst:training}]
from transformers import AutoModelForCausalLM, AutoTokenizer, TrainingArguments
from peft import LoraConfig, get_peft_model
from trl import SFTTrainer

def train(
    model_name: str,
    dataset_path: str,
    output_dir: str,
    use_lora: bool = False,
) -> None:
    tokenizer = AutoTokenizer.from_pretrained(model_name)
    model = AutoModelForCausalLM.from_pretrained(
        model_name,
        torch_dtype="auto",
        attn_implementation="flash_attention_2",
    )
    if use_lora:
        lora_cfg = LoraConfig(
            r=16, lora_alpha=32,
            target_modules=["q_proj", "v_proj"],
            lora_dropout=0.05,
        )
        model = get_peft_model(model, lora_cfg)

    training_args = TrainingArguments(
        output_dir=output_dir,
        num_train_epochs=5,
        per_device_train_batch_size=4,
        gradient_accumulation_steps=4,
        learning_rate=2e-5,
        lr_scheduler_type="cosine",
        warmup_ratio=0.05,
        bf16=True,
        logging_steps=50,
        save_strategy="epoch",
        evaluation_strategy="epoch",
    )
    trainer = SFTTrainer(
        model=model,
        tokenizer=tokenizer,
        train_dataset=load_dataset(dataset_path)["train"],
        eval_dataset=load_dataset(dataset_path)["validation"],
        args=training_args,
    )
    trainer.train()
\end{lstlisting}

\begin{lstlisting}[caption={Regex-метрики оценки структуры ответа},
                   label={lst:metrics}]
import re
import json
from dataclasses import dataclass

@dataclass
class StructureMetrics:
    has_relevant_contexts: bool
    has_analysis: bool
    has_label: bool
    label_valid: bool
    json_parseable: bool

    @property
    def score(self) -> float:
        fields = [
            self.has_relevant_contexts,
            self.has_analysis,
            self.has_label,
            self.label_valid,
            self.json_parseable,
        ]
        return sum(fields) / len(fields)

def evaluate_structure(response: str) -> StructureMetrics:
    """Оценивает структурное соответствие ответа ожидаемому формату."""
    try:
        data = json.loads(response)
        parseable = True
    except json.JSONDecodeError:
        return StructureMetrics(False, False, False, False, False)

    label = data.get("label")
    return StructureMetrics(
        has_relevant_contexts=isinstance(data.get("relevant_contexts"), list),
        has_analysis=bool(data.get("analysis", "").strip()),
        has_label="label" in data,
        label_valid=label in (0, 1, 2),
        json_parseable=parseable,
    )
\end{lstlisting}

\endgroup

\chapter{Скриншоты интерфейса платформы разметки}\label{app:ui}

\begin{figure}[H]
\centering
\includegraphics[width=\textwidth,height=0.62\textheight,keepaspectratio]{code/interface_photo.png}
\caption{Диалоговое окно ассистента разметчика (скриншот интерфейса)}\label{fig:ui-screenshot}
\end{figure}

На рисунке~\ref{fig:ui-screenshot} приведён скриншот диалогового окна
инструмента анализа информационных контекстов в интерфейсе платформы
текстовой разметки (ООО «Яндекс.Технологии»): видна форма запроса,
область информационных контекстов и структурированный ответ ассистента.

\end{document}
